
\section{Second-order character sums: $O(p^{-1})$ where Weil gives $O(p^{-1/2})$}\label{sec:dl}

\paragraph{What this section proves and how to read it.}
A bound is \emph{second order} if it is $O(r^2)=O(1/p)$ where one
Weil bound gives $O(r)$. Three such bounds are proved, for a
four-cycle kernel, a two-variable sum, and trace functions with one
free label (called \emph{moving pinned values} below; ``root'' in
lemma names refers to pinned values, not to the Part~V principal
law). A reader willing to accept \cref{thm:dl:K,thm:dl:H2} as stated
can skip the sheaf arguments; external inputs are collected in
\cref{app:external}.
Weil's bound gives square-root cancellation in a sum of one quadratic
character of one polynomial; after the normalizations of the local
laws this square-root cancellation is a relative error of order
$r=p^{-1/2}$. The spectral criteria of
\cref{sec:spectral}, as they are applied in Part~V, are
conditional: part of a labelling is fixed, and the joint law of the remaining
labels is compared with a product law. Errors of two one-variable counts can
then correlate, and several terms must be of order $r^2$, uniformly in the
fixed labels. This section proves the three second-order estimates the
construction needs, in the generality in which they are used. No graph occurs
in \cref{subsec:dl:K,subsec:dl:H2}, and \cref{subsec:moving-root} uses a half
law only through its one-vertex conditional laws, except for the density bound
of \cref{lem:wt:stars}(a), which uses the tuple densities of a graph law.
\begin{enumerate}[label=(\arabic*)]
\item \emph{The four-cycle kernel} $\chi\big(\prod_{i,j\le2}(y_j-x_i)\big)$ on
ordered pairs of distinct labels has norm at most $2/p=2r^2$ for uniform
probability (\cref{thm:dl:K,cor:dl:normalized}), where Cauchy--Schwarz gives
only $r$. This is the norm of a decided complete rectangle (one all of whose pairs are present or absent; \cref{subsec:extraction:partition}) in the witness
kernels (\cref{lem:matching:kernel}) and of the $K_{2,2}$ core of the
completion variance (\cref{lem:push:localvariance}). Split along the diagonal
of two auxiliary labels, it bounds the remainder of the raw pair maps by order
$q=1/p$ (\cref{cor:raw-pair-maps}): this bound is used for the
positive roots of the late middle pieces (\cref{sec:positive-roots}).
\item \emph{The two-variable sum} $\sum_{y,z}\chi\big(f(y)g(z)(y-z)\big)$ with
split squarefree $f,g$ is $O(p)$ rather than $O(p^{3/2})$ (\cref{thm:dl:H2}).
It bounds the complete conditional means of the product rows
(\cref{subsec:advance:product-rows}) and the pair term of the completion
variance. It also bounds the conditional pivot means of weighted star
marginals (\cref{lem:wt:stars}), which bound the higher-star marginals ($|T|\ge3$ in \cref{lem:wt:stars}) of the
dense comparisons and of the six-query mixture of the upper range
(\cref{prop:dense-comparison,prop:le-mixture}).
\item \emph{The moving-root traces} $\tau_S$ describe the integrated-centre
stars of a cut kernel once all labels but one are fixed. They are almost
orthogonal for distinct root sets, also against character weights
(\cref{lem:moving-root-moments}). Ordered pivot groups turn this into the
star marginal $M_\star$ and the star mean $z_\star$ of the upper root
(\cref{prop:advance:ordered-star,lem:star-mean}), which enter
\cref{lem:record:direct-star-inputs,lem:upper-completed-root}; the uniform
forms of the moments enter the conditional row bounds of the product rows
(\cref{cor:moving-root-uniform,cor:moving-root-conditional}).
\end{enumerate}
The first estimate is representation theory of $\mathrm{PGL}_2(\Fp)$ together
with Jacobi sums. The second and third come from Deligne's theorem applied to
the Katz--Sarnak sheaves of the family of curves $y^2=f(x)(x-z)$; the external
inputs are stated in \cref{app:external}. One further second-order input is
not proved here: the twisted second moment
$\sum_\lambda\chi(\lambda)T_\lambda^2$ of the Legendre traces
$T_\lambda=\sum_y\chi\big(y(y-1)(y-\lambda)\big)$, to which the
graph moment $M_{G_5}$ behind the variance savings (the graph
$G_5$ and its moment are defined in
\cref{lem:record:weighted-moment}) reduces, is evaluated through
a cusp form of
weight four (\cref{thm:ahlgren-ono}) and enters
\cref{lem:record:weighted-moment} directly.
\Cref{subsec:second-order-uses} lists where each estimate enters.
Throughout, $p$ is an odd prime, $r=p^{-1/2}$, $\chi=\chi_p$ is the quadratic
character with $\chi(0)=0$, and $\E_x$ is the average over $x\in\Fp$.

\subsection{The four-cycle kernel}\label{subsec:dl:K}

For four distinct points, the character kernel of a complete
rectangle $K_{2,2}$ is the quadratic character of a cross-ratio
(it is $0$ when points coincide), so it is invariant under
$\mathrm{PGL}_2(\Fp)$
acting on pairs of points of the projective line. The permutation
representation on ordered pairs of distinct points is induced from the split
torus and has multiplicities at most two, so the kernel acts by computable
scalars on all but one isotypic component. The scalars are twisted sums over
the split and the nonsplit torus: products of two Jacobi sums and squares of
sums of modulus $\sqrt p$. This gives spectral norm at most $2p$
for the character matrix (\cref{thm:dl:K}), that is $2/p$ after
normalisation by counting measure
(\cref{cor:dl:normalized}).

\paragraph{Notation.} Points of $\mathbb P^1=\mathbb P^1(\Fp)$ are represented
by the column vectors $(u,1)^{\mathsf T}$, $u\in\Fp$, and
$\infty=(1,0)^{\mathsf T}$; for column vectors $a=(a_0,a_1)^{\mathsf T}$,
$b=(b_0,b_1)^{\mathsf T}$ put $a\wedge b=a_0b_1-a_1b_0$. Let
$\mathcal X=\{(x_1,x_2)\in\mathbb P^1\times\mathbb P^1:x_1\ne x_2\}$, so
$|\mathcal X|=p(p+1)$, and for $x,z\in\mathcal X$ let
\begin{equation}\label{eq:dl:M}
  M(x,z)=\chi\Big(\prod_{i,j\in\{1,2\}}(z_i\wedge x_j)\Big),
\end{equation}
computed with the chosen representatives. For finite points
$z_i\wedge x_j=z_i-x_j$, so on pairs of finite points
$M(x,z)=\chi\big((z_1-x_1)(z_1-x_2)(z_2-x_1)(z_2-x_2)\big)$. We also write $M$
for the matrix $(M(x,z))_{x,z\in\mathcal X}$ and for the operator
$(Mh)(x)=\sum_{z\in\mathcal X}M(x,z)h(z)$ on $\C^{\mathcal X}$ with the
standard inner product $\ip hk=\sum_x\overline{h(x)}k(x)$.

If $\{z_1,z_2\}\cap\{x_1,x_2\}=\emptyset$, the product in \eqref{eq:dl:M} is
the cross-ratio
$[z_1,z_2;x_1,x_2]=\frac{(z_1\wedge x_1)(z_2\wedge x_2)}{(z_1\wedge x_2)(z_2\wedge x_1)}$
times the nonzero square $\big((z_1\wedge x_2)(z_2\wedge x_1)\big)^2$, so
$M(x,z)=\chi([z_1,z_2;x_1,x_2])$; otherwise $M(x,z)=0$. The cross-ratio does
not depend on the representatives and is invariant under
$\mathrm{GL}_2(\Fp)$. Thus $M(x,z)$ depends only on the orbit of the
configuration $(x,z)$, which is why its spectrum can be computed by
representation theory.

Let $\Gamma=\mathrm{PGL}_2(\Fp)$, of order $p(p^2-1)$, acting on $\mathbb P^1$
by matrices on column vectors and diagonally on $\mathcal X$, and let
$\Lambda(g)h(y)=h(g^{-1}y)$ be the permutation representation on
$\C^{\mathcal X}$. We fix a nonsquare $D\in\Fp^\times$ and an element
$\sqrt D\in\F_{p^2}$ with $(\sqrt D)^2=D$, so $\F_{p^2}=\Fp(\sqrt D)$; for
$\zeta\in\F_{p^2}$, $\bar\zeta=\zeta^p$, $N\zeta=\zeta\bar\zeta\in\Fp$ and
$\tr\zeta=\zeta+\bar\zeta\in\Fp$. Let $\mathsf U=\{s\in\F_{p^2}:Ns=1\}$, a
cyclic group of order $p+1$ (the kernel of the norm, which maps
$\F_{p^2}^\times$ onto $\Fp^\times$), and $\eta$ its unique character of order
$2$; for $s\in\mathsf U$, $s^{-1}=\bar s$ and $\tr s=s+s^{-1}$. The map
$\zeta\mapsto\zeta/\bar\zeta=\zeta^{1-p}$ is a homomorphism
$\F_{p^2}^\times\to\mathsf U$ with kernel $\Fp^\times$, hence onto (orders
$p^2-1$ and $p+1$); the elements $\zeta=m+\sqrt D$, $m\in\Fp$, together with
$\zeta=1$ represent the cosets of $\Fp^\times$, and $\zeta=1$ is the one with
$\zeta/\bar\zeta=1$. In $\Gamma$ we use the elements
$u=\begin{psmallmatrix}1&1\\0&1\end{psmallmatrix}$,
$h_a=\begin{psmallmatrix}a&0\\0&1\end{psmallmatrix}$ ($a\in\Fp^\times$), and,
for $\xi\in\F_{p^2}\setminus\Fp$, the image $g_\xi$ in $\Gamma$ of the
$\Fp$-linear map $\zeta\mapsto\xi\zeta$ of $\F_{p^2}=\Fp\oplus\Fp\sqrt D$ (in
the basis $1,\sqrt D$); $g_\xi$ has the eigenvalues $\xi,\bar\xi$.
Multiplicative characters vanish at $0$.

\begin{lemma}[Invariance]\label{lem:dl:invariance}
\begin{enumerate}[label=\textup{(\alph*)}]
\item The value $M(x,z)$ does not depend on the chosen representatives.
  Moreover $M(x,z)=M(z,x)$, and $M(gx,gz)=M(x,z)$ for every
  $g\in\mathrm{GL}_2(\Fp)$. Hence $M$ is real symmetric and commutes with
  $\Lambda(g)$ for every $g\in\Gamma$.
\item $\Gamma$ acts transitively on $\mathcal X$, and the stabilizer of
  $x_0=(0,\infty)$ is the image $\mathsf T=\{h_a:a\in\Fp^\times\}$ of the
  diagonal matrices, of order $p-1$. Hence
  $\C^{\mathcal X}\cong\mathrm{Ind}_{\mathsf T}^\Gamma\mathbf 1$.
\end{enumerate}
\end{lemma}

\begin{proof}
(a) For column vectors and $g\in\mathrm{GL}_2(\Fp)$,
$(ga)\wedge(gb)=\det g\,(a\wedge b)$, and $(ca)\wedge b=c(a\wedge b)=a\wedge(cb)$.
Each of the four points $x_1,x_2,z_1,z_2$ occurs in exactly two of the four
factors $z_i\wedge x_j$. Replacing a representative $a$ by $ca$
($c\in\Fp^\times$) therefore multiplies the product by $c^2$, and replacing
every point by its image under $g$ multiplies it by $(\det g)^4$ times a
product of squares of the normalizing constants of the images; in all cases
the product is multiplied by a nonzero square, which does not change $\chi$.
Moreover
$\prod_{i,j}(x_i\wedge z_j)=\prod_{i,j}(-(z_j\wedge x_i))=\prod_{i,j}(z_i\wedge x_j)$,
so $M$ is symmetric; it is real. If $g\in\Gamma$, then
$(M\Lambda(g)h)(x)=\sum_zM(x,z)h(g^{-1}z)=\sum_zM(x,gz)h(z)
=\sum_zM(g^{-1}x,z)h(z)=(\Lambda(g)Mh)(x)$.
(b) Given $(x_1,x_2)\in\mathcal X$ with representatives $a_1,a_2$, the matrix
with columns $a_1,a_2$ is invertible and maps $\infty=(1,0)^{\mathsf T}$ to
$x_1$ and $0=(0,1)^{\mathsf T}$ to $x_2$; so $\Gamma$ is transitive on
$\mathcal X$. A matrix fixing the points $0$ and $\infty$ has both standard
basis vectors as eigenvectors, so it is diagonal, and its image in $\Gamma$ is
some $h_a$. Since $|\Gamma|/|\mathsf T|=p(p+1)=|\mathcal X|$, the permutation
representation on $\mathcal X$ is $\mathrm{Ind}_{\mathsf T}^\Gamma\mathbf 1$.
\end{proof}

\begin{lemma}[Classes and characters of $\Gamma$]\label{lem:dl:group}
\begin{enumerate}[label=\textup{(\alph*)}]
\item For $\xi\in\F_{p^2}\setminus\Fp$, the conjugacy class of $g_\xi$ in
  $\Gamma$ depends only on $s=\xi/\bar\xi\in\mathsf U\setminus\{1\}$; we write
  $g_s$ for it. Moreover $h_a$ is conjugate to $h_{1/a}$ and $g_s$ to
  $g_{1/s}$, and for every class function $F$ on $\Gamma$
  \begin{equation}\label{eq:dl:classsum}
    \sum_{g\in\Gamma}F(g)=F(1)+(p^2-1)F(u)
    +\frac{p(p+1)}2\sum_{a\in\Fp^\times\setminus\{1\}}F(h_a)
    +\frac{p(p-1)}2\sum_{s\in\mathsf U\setminus\{1\}}F(g_s).
  \end{equation}
\item The irreducible complex representations of $\Gamma$ are, up to
  isomorphism, the following, with the characters shown at $u$, $h_a$
  ($a\ne1$) and $g_s$ ($s\ne1$):
  \begin{center}\small
  \begin{tabular}{@{}lllllr@{}}
  \toprule
  representation & dimension & $u$ & $h_a$ & $g_s$ & number\\
  \midrule
  $\mathbf 1$ & $1$ & $1$ & $1$ & $1$ & $1$\\
  $\mathrm{sgn}$ & $1$ & $1$ & $\chi(a)$ & $\eta(s)$ & $1$\\
  $\mathrm{St}$ & $p$ & $0$ & $1$ & $-1$ & $1$\\
  $\mathrm{St}\otimes\mathrm{sgn}$ & $p$ & $0$ & $\chi(a)$ & $-\eta(s)$ & $1$\\
  $\pi(\mu)$, $\mu\in\widehat{\Fp^\times}$, $\mu^2\ne1$ & $p+1$ & $1$
    & $\mu(a)+\bar\mu(a)$ & $0$ & $(p-3)/2$\\
  $D_\nu$, $\nu\in\widehat{\mathsf U}$, $\nu^2\ne1$ & $p-1$ & $-1$ & $0$
    & $-(\nu(s)+\bar\nu(s))$ & $(p-1)/2$\\
  \bottomrule
  \end{tabular}
  \end{center}
  Here $\pi(\mu)\cong\pi(\bar\mu)$ and $D_\nu\cong D_{\bar\nu}$, and the last
  column is the number of isomorphism classes of each kind; all characters are
  real. $\mathrm{St}$ is the complement of $\mathbf 1$ in the permutation
  representation of $\Gamma$ on $\mathbb P^1$.
\end{enumerate}
\end{lemma}

\begin{proof}
We use \cref{thm:prelim:gl2} with $q=p$. Scaling a matrix does not change its
type, so every $g\ne1$ in $\Gamma$ has a well-defined type (b), (c) or (d).

(a) Let $\xi,\xi'\in\F_{p^2}\setminus\Fp$ with
$\xi'/\bar\xi'\in\{\xi/\bar\xi,\bar\xi/\xi\}$. The multiplications by $\xi'$
and by $\bar\xi'$ have the same eigenvalues $\{\xi',\bar\xi'\}$, so they are
conjugate in $\mathrm{GL}_2(\Fp)$ (\cref{thm:prelim:gl2}; explicitly, the
$\Fp$-linear map $\zeta\mapsto\bar\zeta$ conjugates one into the other).
Replacing $\xi'$ by $\bar\xi'$ if necessary, we may therefore assume
$\xi'/\bar\xi'=\xi/\bar\xi$. Then $c=\xi'/\xi$ satisfies $\bar c=c$, so
$c\in\Fp^\times$, and the multiplication by $\xi'=c\xi$ is $c$ times the
multiplication by $\xi$; hence $g_{\xi'}=g_\xi$ in $\Gamma$. This shows that
the class depends only on $s$ and that $g_s\sim g_{1/s}$. Similarly $h_{1/a}$
is the image of $\mathrm{diag}(1,a)$, which is conjugate to $\mathrm{diag}(a,1)$.

For \eqref{eq:dl:classsum} we count the elements of $\Gamma$ of each type; each
element of $\Gamma$ has exactly $p-1$ preimages in $\mathrm{GL}_2(\Fp)$. By
\cref{thm:prelim:gl2}(iv), a conjugacy class of $\mathrm{GL}_2(\Fp)$ of type
(b), (c) or (d) has $p^2-1$, $p(p+1)$ or $p(p-1)$ elements.
\emph{Type (b).} For $x\in\Fp^\times$ the non-scalar matrices with
characteristic polynomial $(T-x)^2$ form one class, so there are
$(p-1)(p^2-1)$ such matrices and $p^2-1$ elements of type (b) in $\Gamma$. All
of them are conjugate to $u$: the matrix
$\begin{psmallmatrix}x&1\\0&x\end{psmallmatrix}$ represents the same element as
$\begin{psmallmatrix}1&1/x\\0&1\end{psmallmatrix}
=\mathrm{diag}(1/x,1)\,u\,\mathrm{diag}(1/x,1)^{-1}$.
\emph{Type (c).} For $a\in\Fp^\times\setminus\{1\}$ let $C_a\subseteq\Gamma$ be
the set of elements with a preimage whose eigenvalues are $x,ax$ for some $x$;
then $C_a=C_{1/a}$, every element of $C_a$ is conjugate to $h_a$, and the $C_a$
partition the elements of type (c). The eigenvalue sets $\{x,ax\}$,
$x\in\Fp^\times$, are $p-1$ distinct sets if $a\ne-1$ and $(p-1)/2$ if $a=-1$.
So $|C_a|=p(p+1)$ for $a\ne\pm1$ and $|C_{-1}|=p(p+1)/2$. In the sum
$\frac{p(p+1)}2\sum_{a\ne1}F(h_a)$ each set $C_a$ with $a\ne-1$ appears twice
(as $a$ and $1/a$) with weight $\frac{p(p+1)}2$, and $C_{-1}$ once; so this sum
is $\sum_{g\text{ of type (c)}}F(g)$.
\emph{Type (d).} For $s\in\mathsf U\setminus\{1\}$ let $C'_s\subseteq\Gamma$ be
the set of elements with a preimage whose eigenvalues $\xi,\bar\xi$ satisfy
$\xi/\bar\xi\in\{s,1/s\}$; by \cref{thm:prelim:gl2} and the first part of
(a), every element of $C'_s$ is conjugate to $g_s$. The equation
$\xi/\bar\xi=s$ has $p-1$ solutions, and the eigenvalue sets $\{\xi,\bar\xi\}$
with $\xi/\bar\xi=s$ are $p-1$ distinct sets if $s\ne-1$ (since then
$\bar\xi/\xi=1/s\ne s$) and $(p-1)/2$ if $s=-1$. So $|C'_s|=p(p-1)$ for
$s\ne\pm1$ and $|C'_{-1}|=p(p-1)/2$, and as for type (c) the last sum in
\eqref{eq:dl:classsum} is $\sum_{g\text{ of type (d)}}F(g)$. (The count
$1+(p^2-1)+\frac{p(p+1)}2(p-2)+\frac{p(p-1)}2p=p^3-p$ confirms that every
element is counted once.)

(b) The irreducible representations of $\Gamma$ are those of
$\mathrm{GL}_2(\Fp)$ on which the scalar matrices act trivially. By
\cref{thm:prelim:gl2}(iii) and Schur's lemma (a scalar matrix acts on an
irreducible representation by its character value divided by the dimension),
the scalar $xI$ acts on $U_\alpha$ and on $V_\alpha$ by $\alpha(x)^2$, on
$W_{\alpha,\beta}$ by $\alpha\beta(x)$ and on $X_\varphi$ by $\varphi(x)$.
Hence the representations of $\Gamma$ are: $U_\alpha$, $V_\alpha$ with
$\alpha^2=1$, that is $\alpha\in\{1,\chi\}$ (we put $\mathbf 1=U_1$,
$\mathrm{sgn}=U_\chi$, $\mathrm{St}=V_1$, $\mathrm{St}\otimes\mathrm{sgn}=V_\chi$);
$W_{\mu,\bar\mu}=:\pi(\mu)$ with $\mu\ne\bar\mu$, that is $\mu^2\ne1$; and
$X_\varphi$ with $\varphi|_{\Fp^\times}=1$ and $\varphi^p\ne\varphi$. A
character $\varphi$ of $\F_{p^2}^\times$ trivial on $\Fp^\times$ factors through
the surjection $\zeta\mapsto\zeta/\bar\zeta$ onto $\mathsf U$:
$\varphi(\zeta)=\nu(\zeta/\bar\zeta)$ for a unique $\nu\in\widehat{\mathsf U}$;
since $(\zeta/\bar\zeta)^p=\bar\zeta/\zeta$, $\varphi^p=\varphi$ iff
$\bar\nu=\nu$ iff $\nu^2=1$; we put $D_\nu=X_\varphi$. The isomorphisms are
those of \cref{thm:prelim:gl2}(ii), and the numbers follow ($p-3$ characters
$\mu$ of $\Fp^\times$ and $p-1$ characters $\nu$ of $\mathsf U$ have square
$\ne1$). The values: $u$ is of type (b) with $x=1$, $h_a$ of type (c) with
$(x,y)=(a,1)$, and $g_\xi$ of type (d) with $\zeta=\xi$; so
$\varphi(\xi)+\varphi(\bar\xi)=\nu(s)+\nu(\bar s)=\nu(s)+\bar\nu(s)$. Finally
$\xi\mapsto\chi(N\xi)$ is a character of $\F_{p^2}^\times$ that is trivial on
$\Fp^\times$ ($Nc=c^2$) and nontrivial (the norm is onto $\Fp^\times$), so it is
$\nu_0(\xi/\bar\xi)$ with $\nu_0$ of order $2$, that is $\chi(N\xi)=\eta(s)$;
this gives the values of $\mathrm{sgn}$ and $\mathrm{St}\otimes\mathrm{sgn}$ at
$g_s$. All values are real.
\end{proof}

\begin{lemma}[Decomposition]\label{lem:dl:decomp}
As a representation of $\Gamma$,
\[
  \C^{\mathcal X}\cong\mathbf 1\oplus\mathrm{St}^{\oplus2}
  \oplus(\mathrm{St}\otimes\mathrm{sgn})
  \oplus\bigoplus_{\{\mu,\bar\mu\}}\pi(\mu)\oplus\bigoplus_{\{\nu,\bar\nu\}}D_\nu,
\]
the sums over the unordered pairs with $\mu^2\ne1$, resp.\ $\nu^2\ne1$;
$\mathrm{sgn}$ does not occur.
\end{lemma}

\begin{proof}
By \cref{lem:dl:invariance}(b) and Frobenius reciprocity, the multiplicity of
an irreducible $\pi$ is
$m_\pi=\dim\pi^{\mathsf T}=\frac1{p-1}\sum_{a\in\Fp^\times}\chi_\pi(h_a)
=\frac1{p-1}\big(\dim\pi+\sum_{a\ne1}\chi_\pi(h_a)\big)$. With the table of
\cref{lem:dl:group}(b) and $\sum_{a\ne1}\chi(a)=-1$,
$\sum_{a\ne1}(\mu+\bar\mu)(a)=-2$ ($\mu\ne1$): $m_{\mathbf 1}=1$,
$m_{\mathrm{sgn}}=0$, $m_{\mathrm{St}}=\frac{p+p-2}{p-1}=2$,
$m_{\mathrm{St}\otimes\mathrm{sgn}}=1$, $m_{\pi(\mu)}=\frac{p+1-2}{p-1}=1$,
$m_{D_\nu}=1$. (Check:
$1+2p+p+\frac{p-3}2(p+1)+\frac{p-1}2(p-1)=p(p+1)$.)
\end{proof}

\begin{lemma}[Eigenvalue on a multiplicity-one constituent]\label{lem:dl:schur}
Let $\pi$ be an irreducible representation of $\Gamma$ of multiplicity one in
$\C^{\mathcal X}$, and $V_\pi$ its isotypic subspace. Then $M$ acts on $V_\pi$
as the scalar
\[
  \lambda_\pi=\frac1{|\Gamma|}\sum_{g\in\Gamma}\chi_\pi(g)\,\Phi(g),\qquad
  \Phi(g)=\sum_{x\in\mathcal X}M(x,gx).
\]
$\Phi$ is a class function.
\end{lemma}

\begin{proof}
$M$ commutes with $\Lambda(\Gamma)$ (\cref{lem:dl:invariance}(a)), so it
preserves $V_\pi\cong\pi$ and acts there by a scalar $\lambda_\pi$ (Schur's
lemma). $P_\pi=\frac{\dim\pi}{|\Gamma|}\sum_g\overline{\chi_\pi(g)}\Lambda(g)$ is
the orthogonal projection onto $V_\pi$, so
$\lambda_\pi\dim\pi=\mathrm{Tr}(MP_\pi)
=\frac{\dim\pi}{|\Gamma|}\sum_g\overline{\chi_\pi(g)}\,\mathrm{Tr}(M\Lambda(g))$.
The matrix of $\Lambda(g)$ has the entries $[y=gx]$, so
$\mathrm{Tr}(M\Lambda(g))=\sum_xM(x,gx)=\Phi(g)$; the characters are real
(\cref{lem:dl:group}(b)). $\Phi(kgk^{-1})=\Phi(g)$ by the substitution
$x\mapsto kx$ and \cref{lem:dl:invariance}(a).
\end{proof}

For $a\in\Fp^\times$ put $\Upsilon(a)=\sum_{t\in\Fp}\chi\big(t(a-t)(at-1)\big)$,
and for $s\in\mathsf U$ put
$R(s)=\sum_{t\in\mathsf U\setminus\{1\}}\chi(\tr s-\tr t)$.

\begin{lemma}[The function $\Phi$]\label{lem:dl:Phi}
\begin{enumerate}[label=\textup{(\roman*)}]
\item $\Phi(1)=0$ and $\Phi(u)=-p(1+\chi(-1))$.
\item $\Phi(h_a)=(p-1)(\Upsilon(a)-1)$ for $a\ne1$; $\Upsilon(1/a)=\Upsilon(a)$
  and $\Upsilon(1)=-\chi(-1)$.
\item $\Phi(g_s)=-(p+1)\,\eta(s)\,R(s)$ for $s\in\mathsf U\setminus\{1\}$.
\end{enumerate}
\end{lemma}

\begin{proof}
(i) $x_1\wedge x_1=0$, so $M(x,x)=0$. The element $u$ maps $(v,1)^{\mathsf T}$
to $(v+1,1)^{\mathsf T}$ and fixes $\infty$. If $x$ contains $\infty$, the
factor $(u\infty)\wedge\infty$ vanishes. For finite $x$ with $d=x_1-x_2\ne0$
the four factors $(ux)_i-x_j$ are $1$, $1+d$, $1-d$, $1$, so
$\Phi(u)=p\sum_{d\ne0}\chi(1-d^2)$. Now
$\sum_{d\in\Fp}\chi(1-d^2)=\chi(-1)\sum_d\chi((d-1)(d+1))=-\chi(-1)$ by the
two-root identity \eqref{eq:two-root}, and the term $d=0$ is $1$; so
$\Phi(u)=p(-\chi(-1)-1)$.
(ii) $h_a$ fixes $0$ and $\infty$, so pairs containing one of them contribute
$0$. For $x_1,x_2\in\Fp^\times$ write $x_2=tx_1$, $t\notin\{0,1\}$; the four
factors are $(a-1)x_1$, $(a-t)x_1$, $(at-1)x_1$, $(a-1)tx_1$, with product
$(a-1)^2x_1^4\,t(a-t)(at-1)$. Hence
$\Phi(h_a)=(p-1)\sum_{t\ne0,1}\chi(t(a-t)(at-1))=(p-1)(\Upsilon(a)-1)$: the
term $t=0$ of $\Upsilon(a)$ is $0$ and the term $t=1$ is $\chi((a-1)^2)=1$.
Next, $\Upsilon(1/a)=\sum_t\chi\big(t(1-at)(t-a)a^{-2}\big)=\Upsilon(a)$, and
$\Upsilon(1)=\sum_t\chi(-t(t-1)^2)=\sum_{t\ne1}\chi(-t)=-\chi(-1)$.
(iii) Identify $\Fp^2$ with $\F_{p^2}$ by $(a_0,a_1)\mapsto\alpha=a_0+a_1\sqrt D$.
For $\alpha$, $\beta$ corresponding to $a$, $b$ one checks
$\bar\alpha\beta-\alpha\bar\beta=2\sqrt D(a_0b_1-a_1b_0)$, so
$a\wedge b=(\bar\alpha\beta-\alpha\bar\beta)/(2\sqrt D)$. The points of
$\mathbb P^1$ are the cosets $\F_{p^2}^\times/\Fp^\times$, and $g_\xi$ acts by
multiplication by $\xi$. The multiplications by elements of $\F_{p^2}^\times$
form a subgroup of $\mathrm{GL}_2(\Fp)$ whose image $\mathsf T'\subseteq\Gamma$
acts simply transitively on $\mathbb P^1$ and commutes with $g_\xi$. For
$k\in\mathsf T'$, $M(kx,g_\xi kx)=M(kx,kg_\xi x)=M(x,g_\xi x)$ by
\cref{lem:dl:invariance}(a); every $x\in\mathcal X$ is mapped by exactly one
$k\in\mathsf T'$ to a pair with first point $[1]$, so
$\Phi(g_\xi)=(p+1)\sum_wM\big(([1],[w]),([\xi],[\xi w])\big)$, the sum over the
$p$ cosets $[w]\ne[1]$. The numerators $2\sqrt D(z_i\wedge x_j)$ of the four
factors are $\bar\xi-\xi$, $\bar\xi w-\xi\bar w$, $\bar\xi\bar w-\xi w$ and
$N(w)(\bar\xi-\xi)$, and the common denominator $(2\sqrt D)^4=16D^2$ is a
square. Write $\xi=\xi_0+\xi_1\sqrt D$ with $\xi_1\ne0$; then
$(\bar\xi-\xi)^2=4\xi_1^2D$ is a nonsquare. Further
\[
  (\bar\xi w-\xi\bar w)(\bar\xi\bar w-\xi w)=N(w)\tr(\xi^2)-N(\xi)\tr(w^2)
  =N(w)N(\xi)\big(\tr s-\tr t\big),\qquad s=\xi/\bar\xi,\ t=w/\bar w,
\]
since $\tr(\zeta^2)/N\zeta=\zeta/\bar\zeta+\bar\zeta/\zeta$. Hence the summand
is $-\chi(N\xi)\chi(\tr s-\tr t)=-\eta(s)\chi(\tr s-\tr t)$
(\cref{lem:dl:group}(b) for $\chi(N\xi)=\eta(s)$), and as $[w]$ runs over the
cosets $\ne[1]$, $t=w/\bar w$ runs over $\mathsf U\setminus\{1\}$ bijectively.
\end{proof}

We also need two identities. Let $s\in\mathsf U\setminus\{1\}$ and write
$s=w/\bar w$ with $w=m+\sqrt D$. Then $\tr s-2=-(s-1)(\bar s-1)=-N(s-1)$ and
$s-1=(w-\bar w)/\bar w=2\sqrt D/\bar w$, so $N(s-1)=-4D/N(w)$ and
$\chi(\tr s-2)=\chi(D)\chi(Nw)=-\eta(s)$. Since $\chi(2-\tr t)=\chi(-1)\chi(\tr t-2)$ and $\sum_{t\in\mathsf U}\eta(t)=0$,
\begin{equation}\label{eq:dl:trace2}
  \sum_{t\in\mathsf U}\chi(2-\tr t)=\chi(-1),\qquad
  \sum_{s\in\mathsf U}\psi(s)\chi(\tr s-2)=-\sum_{s\ne1}(\psi\eta)(s)=1
  \quad(\psi\in\widehat{\mathsf U},\ \psi\ne\eta).
\end{equation}
For $\psi\in\widehat{\mathsf U}$ and $\mu\in\widehat{\Fp^\times}$ put
\[
  \mathcal N(\psi)=\sum_{s,t\in\mathsf U}\psi(s)\chi(\tr s-\tr t),\qquad
  \mathcal S(\mu)=\sum_{a\in\Fp^\times}\mu(a)\Upsilon(a),
\]
and let $\Jac(\psi_1,\psi_2)=\sum_{v\in\Fp}\psi_1(v)\psi_2(1-v)$ be the Jacobi
sum of \cref{thm:jacobi}. Since $\tr(s^{-1})=\tr s$ and
$\Upsilon(1/a)=\Upsilon(a)$, $\mathcal N(\bar\psi)=\mathcal N(\psi)$ and
$\mathcal S(\bar\mu)=\mathcal S(\mu)$. Since $\Fp^\times\subseteq\F_{p^2}^\times$
is a subgroup of a finite abelian group, $\chi$ extends to a character
$\tilde\chi$ of $\F_{p^2}^\times$; we fix one ($\tilde\chi(0)=0$). For
$\kappa\in\widehat{\mathsf U}$ put
$\mathcal C(\kappa)=\sum_{v\in\mathsf U}\kappa(v)\tilde\chi(v-1)$.

\begin{lemma}[Line sums]\label{lem:dl:line}
Let $\Xi$ be a character of $\F_{p^2}^\times$ whose restriction to
$\Fp^\times$ is nontrivial, and $\alpha\in\F_{p^2}\setminus\Fp$. Then
$\big|\sum_{m\in\Fp}\Xi(m+\alpha)\big|=\sqrt p$.
\end{lemma}

\begin{proof}
$m+\alpha\ne0$ for all $m$, and
$|\sum_m\Xi(m+\alpha)|^2=p+\sum_{m\ne n}\Xi\big(\frac{m+\alpha}{n+\alpha}\big)$.
We claim that $(m,n)\mapsto w=\frac{m+\alpha}{n+\alpha}$ is a bijection from
$\{(m,n)\in\Fp^2:m\ne n\}$ onto $\F_{p^2}\setminus\Fp$. If $w=c\in\Fp$, then
$(1-c)\alpha=cn-m\in\Fp$, which forces $c=1$ and $m=n$; so the map lands in
$\F_{p^2}\setminus\Fp$. Given $w\notin\Fp$, the elements $1,w$ are a basis of
$\F_{p^2}$ over $\Fp$, so $m-wn=(w-1)\alpha$ has exactly one solution
$(m,n)\in\Fp^2$, and $m\ne n$ (otherwise $(w-1)(m+\alpha)=0$, impossible); this
solution is the unique preimage. Hence the last sum is
$\sum_{\F_{p^2}^\times}\Xi-\sum_{\Fp^\times}\Xi=0-0$, because $\Xi$ is
nontrivial on $\Fp^\times$, hence on $\F_{p^2}^\times$.
\end{proof}

\begin{lemma}[Torus sums]\label{lem:dl:split}
\begin{enumerate}[label=\textup{(\alph*)}]
\item Let $C$ be a cyclic group of even order, $-1$ its element of order two,
  $\psi$ a character of $C$ and $F:C\to\C$. If $\psi(-1)=-1$ then
  $\sum_{y\in C}\psi(y)F(y^2)=0$. If $\psi(-1)=1$ then
  $\sum_{y\in C}\psi(y)F(y^2)=\sum_\kappa\sum_{w\in C}\kappa(w)F(w)$, the sum
  over the two characters $\kappa$ of $C$ with $\kappa^2=\psi$.
\item \emph{(Split torus.)} If $\mu(-1)=-1$ then $\mathcal S(\mu)=0$. If
  $\mu(-1)=1$, then
  \[
    \mathcal S(\mu)=\chi(-1)\sum_{\kappa}\Jac(\kappa,\chi)\Jac(\bar\kappa\chi,\chi),
  \]
  the sum over the two characters $\kappa$ of $\Fp^\times$ with
  $\kappa^2=\bar\mu$. If $\mu^2\ne1$, or if $\mu=\chi$, then
  $|\mathcal S(\mu)|\le2p$.
\item \emph{(Nonsplit torus.)} Let $\psi\in\widehat{\mathsf U}$ and
  $\rho=(\psi\bar{\tilde\chi})|_{\mathsf U}$. If $\rho(-1)=-1$ then
  $\mathcal N(\psi)=0$; otherwise
  $\mathcal N(\psi)=\sum_{\kappa}\mathcal C(\kappa)^2$, the sum over the two
  $\kappa\in\widehat{\mathsf U}$ with $\kappa^2=\rho$. Moreover
  $|\mathcal C(\kappa)|=\sqrt p$ for every $\kappa\in\widehat{\mathsf U}$; hence
  $|\mathcal N(\psi)|\le2p$ for every $\psi$.
\end{enumerate}
\end{lemma}

\begin{proof}
(a) If $\psi(-1)=-1$, the substitution $y\mapsto-y$ multiplies the sum by
$\psi(-1)=-1$, so it vanishes. If $\psi(-1)=1$: the characters of $C$ form a
cyclic group of even order, in which the characters with value $1$ at $-1$
form the subgroup of index two, that is the squares; so $\psi=\kappa_0^2$, and
the characters with square $\psi$ are $\kappa_0$ and $\kappa_0\chi_C$, where
$\chi_C$ is the character of order two of $C$. As $y\mapsto y^2$ is two-to-one
onto the subgroup $C^2$ of squares and $\kappa_0+\kappa_0\chi_C=2\kappa_0\one_{C^2}$,
$\sum_y\psi(y)F(y^2)=\sum_y\kappa_0(y^2)F(y^2)=2\sum_{w\in C^2}\kappa_0(w)F(w)
=\sum_{w\in C}(\kappa_0+\kappa_0\chi_C)(w)F(w)$.

(b) The terms $t=0$ of $\Upsilon$ vanish. For $t\ne0$ write $a=tb$; then
$t(a-t)(at-1)=t^2(b-1)(t^2b-1)$. With $c=t^2b$ we have
$\mu(a)=\bar\mu(t)\mu(c)$ and $\chi(b-1)=\chi(c-t^2)$, so
\[
  \mathcal S(\mu)=\sum_{c\ne0}\mu(c)\chi(c-1)\,I(c),\qquad
  I(c)=\sum_{t\ne0}\bar\mu(t)\chi(c-t^2).
\]
Apply (a) with $C=\Fp^\times$, $\psi=\bar\mu$ (so $\psi(-1)=\mu(-1)$) and
$F(y)=\chi(c-y)$: if $\mu(-1)=-1$ then $I(c)=0$, and otherwise
$I(c)=\sum_\kappa\sum_{y\ne0}\kappa(y)\chi(c-y)$. For $c\ne0$ the substitution
$y=cv$ gives $\sum_y\kappa(y)\chi(c-y)=(\kappa\chi)(c)\Jac(\kappa,\chi)$. Hence
$\mathcal S(\mu)=\sum_\kappa\Jac(\kappa,\chi)\sum_{c\ne0}(\mu\kappa\chi)(c)\chi(c-1)$,
and $\mu\kappa=\bar\kappa$ (as $\kappa^2=\bar\mu$), while
$\sum_c(\bar\kappa\chi)(c)\chi(c-1)=\chi(-1)\sum_c(\bar\kappa\chi)(c)\chi(1-c)
=\chi(-1)\Jac(\bar\kappa\chi,\chi)$.
Moduli: if $\mu^2\ne1$ then $\kappa^4\ne1$, so $\kappa\ne1$, $\kappa\ne\chi$;
hence $\kappa$, $\chi$, $\kappa\chi$ and $\bar\kappa\chi$, $\chi$,
$\bar\kappa\chi\cdot\chi=\bar\kappa$ are nontrivial, and \cref{thm:jacobi} gives
$|\Jac|=\sqrt p$ for both Jacobi sums. If $\mu=\chi$ is even,
i.e. $\chi(-1)=1$ ($p\equiv1\pmod4$), $\kappa$ has order $4$, so
$\kappa\chi=\bar\kappa$ and $\bar\kappa\chi=\kappa$ are
nontrivial, and again $|\Jac|=\sqrt p$. In both cases
$|\mathcal S(\mu)|\le2p$; if $\mu=\chi$ is odd
($\chi(-1)=-1$), $\mathcal S(\chi)=0$.

(c) For $s,t\in\mathsf U$, $\tr s-\tr t=(s-t)(st-1)/(st)\in\Fp$. Since
$\tilde\chi$ restricts to $\chi$ on $\Fp^\times$ and is multiplicative,
$\chi(\tr s-\tr t)=\tilde\chi(s-t)\tilde\chi(st-1)\bar{\tilde\chi}(st)$ (both
sides vanish iff $t\in\{s,s^{-1}\}$). Put $t=sw$, $w\in\mathsf U$: the summand
$\psi(s)\chi(\tr s-\tr t)$ becomes
$\rho(s)\cdot\bar{\tilde\chi}(w)\tilde\chi(1-w)\cdot\tilde\chi(s^2w-1)$. So
$\mathcal N(\psi)=\sum_{w}\bar{\tilde\chi}(w)\tilde\chi(1-w)
\sum_{s\in\mathsf U}\rho(s)\tilde\chi(s^2w-1)$. Apply (a) with $C=\mathsf U$
(cyclic of even order $p+1$; note $-1\in\mathsf U$) and
$F(y)=\tilde\chi(yw-1)$: the inner sum vanishes if $\rho(-1)=-1$, and otherwise
it is $\sum_\kappa\sum_{y\in\mathsf U}\kappa(y)\tilde\chi(yw-1)
=\sum_\kappa\bar\kappa(w)\mathcal C(\kappa)$ ($v=yw$). Finally
$\bar{\tilde\chi}(w)\tilde\chi(1-w)=\tilde\chi(w^{-1}-1)$, and
$\sum_w\bar\kappa(w)\tilde\chi(w^{-1}-1)=\mathcal C(\kappa)$ ($w\mapsto w^{-1}$).

For $|\mathcal C(\kappa)|$: the map $z\mapsto z^{p-1}=\bar z/z$ sends
$\F_{p^2}^\times$ onto $\mathsf U$ with kernel $\Fp^\times$, the coset of $1$
goes to $1$ (whose term $\tilde\chi(1-1)=0$ vanishes), and the cosets of
$z=m+\sqrt D$, $m\in\Fp$, go bijectively to $\mathsf U\setminus\{1\}$. For such
$z$, $\bar z/z-1=-2\sqrt D/z$. Hence, with the character
$\Xi(z)=\kappa(z^{p-1})\bar{\tilde\chi}(z)$ of $\F_{p^2}^\times$,
\[
  \mathcal C(\kappa)=\tilde\chi(-2\sqrt D)\sum_{m\in\Fp}\Xi(m+\sqrt D),
\]
and $\Xi$ restricts to $\bar\chi=\chi\ne1$ on $\Fp^\times$, because
$z\mapsto\kappa(z^{p-1})$ is trivial there. \Cref{lem:dl:line} with
$\alpha=\sqrt D$ gives $|\mathcal C(\kappa)|=\sqrt p$.
\end{proof}

\begin{theorem}[Four-cycle operator norm]\label{thm:dl:K}
For every odd prime $p$, every eigenvalue of the real symmetric matrix $M$ of
\eqref{eq:dl:M} lies in $[-2p,2p]$. In particular every principal submatrix of
$M$ (for instance the one indexed by the ordered pairs of distinct elements of
$\Fp$) has spectral norm at most $2p$.
\end{theorem}

\begin{proof}
By \cref{lem:dl:decomp} $\C^{\mathcal X}$ is the orthogonal sum of the isotypic
subspaces, which $M$ preserves (it commutes with $\Lambda(\Gamma)$). We
determine the eigenvalues on each.

\emph{$\pi(\mu)$, $\mu^2\ne1$.} By \cref{lem:dl:schur}, \eqref{eq:dl:classsum},
\cref{lem:dl:Phi} and \cref{lem:dl:group}(b),
\begin{align*}
  \lambda&=\frac1{|\Gamma|}\Big[(p^2-1)\Phi(u)
    +\frac{p(p+1)}2\sum_{a\ne1}(\mu+\bar\mu)(a)(p-1)(\Upsilon(a)-1)\Big]\\
  &=-(1+\chi(-1))+\frac12\sum_{a\ne1}(\mu+\bar\mu)(a)(\Upsilon(a)-1).
\end{align*}
Since $\sum_{a\ne1}\mu(a)=-1$ and $\Upsilon(1)=-\chi(-1)$,
$\sum_{a\ne1}\mu(a)(\Upsilon(a)-1)=\mathcal S(\mu)+\chi(-1)+1$, and likewise for
$\bar\mu$ with $\mathcal S(\bar\mu)=\mathcal S(\mu)$. So
$\lambda=\mathcal S(\mu)$, and $|\lambda|\le2p$ by \cref{lem:dl:split}(b).

\emph{$D_\nu$, $\nu^2\ne1$.} Here $\chi_\pi(u)=-1$, $\chi_\pi(h_a)=0$, so
$\lambda=(1+\chi(-1))+\frac12\sum_{s\ne1}(\nu+\bar\nu)(s)\eta(s)R(s)$. For
$\psi\in\widehat{\mathsf U}$ with $\psi\ne\eta$, removing from $\mathcal N(\psi)$
the row $s=1$ and the column $t=1$ (whose common term $\chi(0)$ is $0$) gives,
by \eqref{eq:dl:trace2}, $\sum_{s\ne1}\psi(s)R(s)=\mathcal N(\psi)-\chi(-1)-1$.
With $\psi=\nu\eta$ and $\psi=\bar\nu\eta$ (both $\ne\eta$ as $\nu\ne1$) and
$\mathcal N(\bar\nu\eta)=\mathcal N(\overline{\nu\eta})=\mathcal N(\nu\eta)$, we
get $\lambda=\mathcal N(\nu\eta)$, and $|\lambda|\le2p$ by
\cref{lem:dl:split}(c).

\emph{$\mathrm{St}\otimes\mathrm{sgn}$.} Here $\chi_\pi(u)=0$,
$\chi_\pi(h_a)=\chi(a)$, $\chi_\pi(g_s)=-\eta(s)$, so
\begin{align*}
  \lambda&=\frac12\sum_{a\ne1}\chi(a)(\Upsilon(a)-1)+\frac12\sum_{s\ne1}R(s)\\
  &=\frac12\big(\mathcal S(\chi)+\chi(-1)+1\big)
  +\frac12\big(\mathcal N(\mathbf 1)-\chi(-1)-1\big)
  =\frac12\big(\mathcal S(\chi)+\mathcal N(\mathbf 1)\big),
\end{align*}
and $|\lambda|\le2p$ by \cref{lem:dl:split}(b),(c).

\emph{The trivial representation.} It is spanned by the constant function, with
eigenvalue the row sum $\sum_zM(x_0,z)$, which is the same for all $x_0$
(transitivity). For $x_0=(0,\infty)$: with the representatives $(z_i,1)$, $(0,1)$, $(1,0)$ fixed above, $z_i\wedge0=z_i$ and $z_i\wedge\infty=-1$
for finite $z_i$, while $\infty\wedge\infty=0$; so $M(x_0,z)=\chi(z_1z_2)$ if
$z_1,z_2\in\Fp$ and $0$ otherwise, and
$\sum_{z_1\ne z_2}\chi(z_1z_2)=(\sum_{v}\chi(v))^2-\sum_v\chi(v^2)=-(p-1)$.

\emph{$\mathrm{St}$ (multiplicity $2$).} Let $C(z,x)=\chi(z\wedge x)$ for
$z,x\in\mathbb P^1$. Then $\sum_{z\in\mathbb P^1}C(z,x)C(z,x')=p\,[x=x']$: for
$x=x'$ there are $p$ points $z\ne x$; for $x\ne x'$ both finite,
$\sum_{z\in\Fp}\chi((z-x)(z-x'))=-1$ by \eqref{eq:two-root}, and the term
$z=\infty$ is $\chi(1)^2=1$; for $x'=\infty$, $x$ finite,
$\sum_{z\in\Fp}\chi(z-x)\chi(-1)=0$ and the term $z=\infty$ vanishes. Now
$M(x,z)=\phi_x(z_1)\phi_x(z_2)$ with $\phi_x(z)=C(z,x_1)C(z,x_2)$, so
$\sum_{z\in\mathbb P^1}\phi_x(z)=0$ and
$\phi_x(z)^2=\one[z\notin\{x_1,x_2\}]$. For $f:\mathbb P^1\to\C$ with
$\sum f=0$ put $F_1(x)=f(x_1)$, $F_2(x)=f(x_2)$. Then
\[
  (MF_1)(x)=\sum_{z_1}f(z_1)\phi_x(z_1)\sum_{z_2\ne z_1}\phi_x(z_2)
  =-\sum_{z_1}f(z_1)\phi_x(z_1)^2=-\sum_{z_1\notin\{x_1,x_2\}}f(z_1)
  =f(x_1)+f(x_2),
\]
and, $M(x,z)$ being symmetric in $(z_1,z_2)$, also $MF_2=F_1+F_2$. The space
$V_1=\{F_1+F_2':\ \sum f=\sum f'=0\}$ (with $F'_2(x)=f'(x_2)$) has dimension
$2p$ (if $f(x_1)+f'(x_2)=0$ for all $x_1\ne x_2$, then $f$ is constant, as
$p+1\ge3$, hence $0$, and then $f'=0$), is $\Gamma$-invariant and is isomorphic
to $\mathrm{St}\oplus\mathrm{St}$ (\cref{lem:dl:group}(b): $\mathrm{St}$ is the
space of functions on $\mathbb P^1$ with sum $0$). Since $m_{\mathrm{St}}=2$,
$V_1$ is the $\mathrm{St}$-isotypic subspace, and $M$ acts on it as the
identity on $\mathrm{St}$ tensored with
$\begin{psmallmatrix}1&1\\1&1\end{psmallmatrix}$, with eigenvalues $2$ and $0$.

So every eigenvalue of $M$ is one of $\mathcal S(\mu)$, $\mathcal N(\nu\eta)$,
$\frac12(\mathcal S(\chi)+\mathcal N(\mathbf 1))$, $-(p-1)$, $2$, $0$, all of
modulus at most $2p$. A principal submatrix of a real symmetric matrix is a
compression, so its norm is at most $\norm M\le2p$.
\end{proof}

Every application of \cref{thm:dl:K} in this paper needs it for normalized counting measure, for
functions that vanish on noninjective tuples, and sometimes against
conditional laws with bounded densities. We record these forms once.

\begin{corollary}[The four-cycle kernel for uniform probability]\label{cor:dl:normalized}
Let $a,b\ge2$, let $\varsigma_{ij}\in\{\pm1\}$ ($i\le a$, $j\le b$), and put
\[
  K(x,y)=\prod_{i\le a,\ j\le b}\chi\big(\varsigma_{ij}(y_j-x_i)\big)
  \qquad(x\in\Fp^a,\ y\in\Fp^b).
\]
\begin{enumerate}[label=\textup{(\alph*)}]
\item If $f:\Fp^a\to\C$ vanishes on the tuples with $x_1=x_2$ and
  $h:\Fp^b\to\C$ on the tuples with $y_1=y_2$ (in particular, if both vanish on
  noninjective tuples), then, for uniform probability on $\Fp^a$ and $\Fp^b$,
  \begin{equation}\label{eq:dl:normalized}
    \Big|\E_{x,y}\overline{f(x)}\,K(x,y)\,h(y)\Big|\le\frac2p\,\norm f_2\norm h_2 .
  \end{equation}
  For $a=b=2$: the operator with kernel $K$ relative to uniform probability on
  $\Fp^2$, compressed to the functions that vanish on the diagonal, has norm
  at most $2/p=2r^2$.
\item Let $\nu_1,\nu_2$ be probability measures on $\Fp^a$, $\Fp^b$ with
  $\nu_1\{x_1=x_2\}=\nu_2\{y_1=y_2\}=0$ and densities relative to uniform
  probability at most $D_1$, $D_2$. Then for $f\in L^2(\nu_1)$, $h\in L^2(\nu_2)$,
  \[
    \Big|\E_{x\sim\nu_1,\,y\sim\nu_2}\overline{f(x)}\,K(x,y)\,h(y)\Big|
    \le\frac2p\sqrt{D_1D_2}\,\norm f_{\nu_1}\norm h_{\nu_2}.
  \]
\end{enumerate}
\end{corollary}

\begin{proof}
Since $\chi(\varsigma(y-x))=\chi(\varsigma)\chi(y-x)$, the signs multiply $K$ by a
constant of modulus one; we take them equal to $1$.

(a) First let $a=b=2$. For finite pairs $K(x,y)=M(x,y)$, and $f,h$ are supported
on the set $\mathcal X_{\rm aff}$ of ordered pairs of distinct elements of
$\Fp$. So the left side of \eqref{eq:dl:normalized} is
$p^{-4}|\mathbf f^*M_{\rm aff}\mathbf h|$, where $M_{\rm aff}$ is the principal
submatrix of $M$ on $\mathcal X_{\rm aff}$ and $\mathbf f$, $\mathbf h$ are the
vectors of values. By \cref{thm:dl:K} this is at most
$p^{-4}\cdot2p\,\norm{\mathbf f}_{\ell^2}\norm{\mathbf h}_{\ell^2}
=p^{-4}\cdot2p\cdot p\norm f_2\cdot p\norm h_2=\frac2p\norm f_2\norm h_2$. For
general $a,b$ fix $x'=(x_3,\dots,x_a)$ and $y'=(y_3,\dots,y_b)$. Then
$K(x,y)=K_{2,2}\big((x_1,x_2),(y_1,y_2)\big)\,\alpha(x_1,x_2)\beta(y_1,y_2)c$,
where $K_{2,2}$ is the kernel for $a=b=2$, $\alpha=\prod_{i\le2<j}\chi(y_j-x_i)$,
$\beta=\prod_{j\le2<i}\chi(y_j-x_i)$ and $c=\prod_{i,j>2}\chi(y_j-x_i)$ have
modulus at most one and depend on $(x',y')$. The sliced functions
$f(\cdot,x')\alpha$ and $h(\cdot,y')\beta$ vanish on the diagonal, so the case
$a=b=2$ bounds the inner average over $(x_1,x_2,y_1,y_2)$ by
$\frac2p\norm{f(\cdot,x')}_2\norm{h(\cdot,y')}_2$. Averaging over the
independent $x'$ and $y'$ and applying Cauchy--Schwarz separately in $x'$ and in
$y'$ gives \eqref{eq:dl:normalized}.

(b) Let $d_i$ be the density of $\nu_i$. The left side is
$|\E_{x,y}\overline{(fd_1)(x)}K(x,y)(hd_2)(y)|$ for uniform probability, and
$fd_1$, $hd_2$ vanish where $x_1=x_2$, resp.\ $y_1=y_2$. By (a) it is at most
$\frac2p\norm{fd_1}_2\norm{hd_2}_2$, and
$\norm{fd_1}_2^2=\E[d_1^2|f|^2]\le D_1\E[d_1|f|^2]=D_1\norm f_{\nu_1}^2$, and
likewise for $h$.
\end{proof}

The positive roots of the late middle pieces (\cref{sec:positive-roots})
multiply a root kernel by a polynomial of degree two in quartic traces
$\E_s\psi_y(s)\psi_{y'}(s)$ of two pairs of labels, built from the \emph{pair
traces} $\psi_y(s)=\chi(s-y_0)\chi(s-y_1)$. The maps defined by pair traces are close to rank one in the sense of \eqref{eq:raw-pair-split} below ($\norm R^2\le c_0q^2$ there): their diagonal part is explicit, and their off-diagonal part
is the four-cycle kernel, of norm $O(1/p)$ by \cref{cor:dl:normalized}.

\begin{corollary}[Raw pair maps]\label{cor:raw-pair-maps}
Let $p$ be an odd prime and $q=1/p$. Unless said otherwise every label is
uniform on $\Fp$, and $L^2(\Fp^k)$ carries the uniform probability. For a pair
$y=(y_0,y_1)$ of labels put $\psi_y(s)=\chi(s-y_0)\chi(s-y_1)$.
\begin{enumerate}[label=\textup{(\alph*)}]
\item The maps $h\mapsto\E_s\chi(s-a)h(s)$ ($h\in L^2(\Fp)$) and $F\mapsto\E_{s,t}\chi(s-a)\chi(t-a)F(s,t)$ ($F\in L^2(\Fp^2)$), into functions of $a\in\Fp$, have norms
  $r$ and $\sqrt{q(1-q)}\le r$. The \emph{raw pair map}
  $h\mapsto\E_s\psi_y(s)h(s)$ has norm $\sqrt{q(1-q)}$ into functions of
  $y\in\Fp^2$, and norm at most $r$ into functions of $y$ uniform on the ordered
  pairs of distinct labels. A tensor product of such maps on independent labels
  has the product of their norms as its norm.
\item Let $y$ be uniform on the ordered pairs of distinct labels, let
  $S\colon L^2(\Fp^2)\to L^2(y)$ be the map
  $(SF)(y)=\E_{s,u}\psi_y(s)\psi_y(u)F(s,u)$, and let $d=\sqrt p\,\one_{s=u}$, a
  unit vector of $L^2(\Fp^2)$. Then
  \begin{equation}\label{eq:raw-pair-split}
    S=r(1-2q)\,\one\ip d{\cdot}+R,\qquad Rd=0,\qquad\norm R^2\le c_0q^2,\qquad
    c_0=\frac4{1-q}+2q .
  \end{equation}
  In detail, the space $L^2(\Fp^2)$ is the orthogonal sum of $\C d$, of the
  functions $Df$, where $(Df)(s,u)=\sqrt p\,\one_{s=u}f(s)$ and
  $f\in L^2_0(\Fp)$, and of the functions that vanish on the diagonal
  $\{s=u\}$. On the second space
  $(RDf)(y)=-qr\big(f(y_0)+f(y_1)\big)$ and $\norm{RDf}\le\sqrt2\,qr\norm f$; on
  the third, $R$ agrees with $S$ and has norm at most $2q/\sqrt{1-q}$. For every
  $y$ the row $(s,u)\mapsto\psi_y(s)\psi_y(u)-r(1-2q)\,d(s,u)$ of $R$ has norm at
  most one in $L^2(\Fp^2)$.
\end{enumerate}
\end{corollary}

\begin{proof}
By \eqref{eq:two-root}, $\E_a\chi(s-a)\chi(s'-a)=\one_{s=s'}-q$ for all $s,s'$.

(a) The first map is the Paley operator $P$ of \cref{lem:square-shift}, of norm
$r$ by part (c) there. Write $R_2$ for the second map and $R_3$ for the raw pair
map into functions of $y\in\Fp^2$. Then
\[
 (R_2R_2^*g)(a)=\E_{a'}\big(\E_s\chi(s-a)\chi(s-a')\big)^2g(a'),\qquad
 (R_3^*R_3h)(s)=\E_{s'}\big(\E_{y_0}\chi(s-y_0)\chi(s'-y_0)\big)^2h(s'),
\]
and both kernels are $(\one_{a=a'}-q)^2=(1-2q)\one_{a=a'}+q^2$ (with $s,s'$ in
place of $a,a'$ in the second). The operator with this kernel is $q(1-2q)$ times
the identity on $L^2_0(\Fp)$ and $q(1-q)$ on the constants, so
$\norm{R_2}=\norm{R_3}=\sqrt{q(1-q)}$. For $y$ uniform on distinct pairs: the
uniform probability on $\Fp^2$ is $1-q$ times the law of $y$ plus $q$ times the
uniform law on the diagonal, and on the diagonal
$\E_{y_0}\psi_{(y_0,y_0)}(s)\psi_{(y_0,y_0)}(s')=\E_{y_0}\one_{y_0\notin\{s,s'\}}
=1-2q+q\one_{s=s'}$; hence
\[
 \E_y\psi_y(s)\psi_y(s')=\frac{(1-2q)\one_{s=s'}+q^2-q(1-2q+q\one_{s=s'})}{1-q}
 =\frac{1-2q-q^2}{1-q}\one_{s=s'}-\frac{q(1-3q)}{1-q},
\]
and the Gram operator of the raw pair map is $q(1-2q-q^2)/(1-q)\le q$ on
$L^2_0(\Fp)$ and $q^2$ on the constants. Finally, the norm of a tensor product of
operators is the product of their norms.

(b) The map $D\colon L^2(\Fp)\to L^2(\Fp^2)$ is an isometry onto the functions
supported on the diagonal, and $D\one=d$; this gives the orthogonal
decomposition. Since $y_0\ne y_1$, $\psi_y(s)^2=\one_{s\notin\{y_0,y_1\}}$, so
\[
 (SDf)(y)=\sqrt p\,\E_{s,u}\one_{s=u}\psi_y(s)\psi_y(u)f(s)
 =r\,\E_s\one_{s\notin\{y_0,y_1\}}f(s)=r\,\E f-qr\big(f(y_0)+f(y_1)\big).
\]
For $f=\one$ this is $Sd=r(1-2q)\one$, and $\ip dd=1$ gives $Rd=0$. For
$f\in L^2_0(\Fp)$, $\ip d{Df}=\E f=0$, so $RDf=SDf=-qr(f(y_0)+f(y_1))$; as $y_0$
and $y_1$ are uniform and $\E f(y_0)\overline{f(y_1)}=-\norm f^2/(p-1)$,
$\norm{RDf}^2=2q^3\big(1-\frac1{p-1}\big)\norm f^2\le2q^3\norm f^2$.

A function $F$ that vanishes on the diagonal has $\ip dF=0$, so $RF=SF$. The
kernel $\psi_y(s)\psi_y(u)=\prod_{i\in\{0,1\}}\prod_{v\in\{s,u\}}\chi(v-y_i)$ is the
kernel $K$ of \cref{cor:dl:normalized} with $a=b=2$, the pair $y$ in the role of
$x$ and all signs $+1$. For $G\in L^2(y)$ let $\tilde G$ be its extension by zero
to $\Fp^2$, so that $\norm{\tilde G}_{L^2(\Fp^2)}=(1-q)^{1/2}\norm G$. By
\cref{cor:dl:normalized}(a), $\tilde G$ and $F$ vanishing on the diagonals,
\[
 \big|\ip G{SF}_{L^2(y)}\big|
 =\frac1{1-q}\Big|\E_{y\in\Fp^2}\overline{\tilde G(y)}\,(SF)(y)\Big|
 \le\frac1{1-q}\cdot\frac2p\,(1-q)^{1/2}\norm G\norm F
 =\frac{2q}{\sqrt{1-q}}\norm G\norm F .
\]
For a general $F=\lambda d+Df+F'$, with $f\in L^2_0(\Fp)$ and $F'$ vanishing on
the diagonal, $RF=RDf+RF'$, and Cauchy--Schwarz gives
$\norm{RF}^2\le\big(2q^3+\frac{4q^2}{1-q}\big)\big(\norm f^2+\norm{F'}^2\big)
\le c_0q^2\norm F^2$.

Finally, the row of $S$ at $y$ is $k_y(s,u)=\psi_y(s)\psi_y(u)$, of modulus at
most one, so $\norm{k_y}\le1$; and $\ip d{k_y}=(Sd)(y)=r(1-2q)$, so the row of
$R$ at $y$, $k_y-r(1-2q)d=k_y-\ip d{k_y}d$, is the orthogonal projection of
$k_y$ onto $d^\perp$.
\end{proof}

\noindent For two pairs $y,y'$ of distinct labels the normalized quartic trace
$H=r^{-1}\E_s\psi_y(s)\psi_{y'}(s)$ (the scores of the positive roots,
\cref{def:positive-root}, are of this form) satisfies $H^2=q^{-1}\ip{k_y}{k_{y'}}$,
where $k_y$ is the row of $S$ at $y$, and so, by (b),
\[
 H^2=(1-2q)^2+q^{-1}\ip{\rho_y}{\rho_{y'}},\qquad\rho_y=k_y-r(1-2q)\,d .
\]
The square of a score thus differs from the constant $(1-2q)^2$ by a Gram kernel
of the rows of $R$, divided by $q$. With the notation of
\cref{lem:swapped-rectangle} (in particular that this Gram kernel
is positive semidefinite), the bound $\norm R^2\le c_0q^2$ turns
this
into an error of order $q$ in the square estimates of the positive roots
(\cref{lem:root-square}).

\begin{remark}[The constant $2$]\label{rem:dl:K-sharp}
The formulas of the proof of \cref{thm:dl:K} reproduce the full spectrum of $M$ for every odd
prime $p\le73$ (floating-point computations of two internal referees, the
second up to $p\le89$; evaluating the formulas for every odd prime $p<2000$
gives $\max|\lambda|<2p$ throughout), and $\max|\lambda|/p$ reaches
$2-4.4\cdot10^{-12}$ at $p=71$. These floating-point computations suggest
that the constant $2$ is optimal; they do not prove optimality (and are
not used). The
exact finite consistency tests of \cref{thm:dl:K} are described in
\cref{rem:cert:deligne-tests}; they are not used in any proof.
\end{remark}

\subsection{The complete three-by-three character kernel}
\label{subsec:cubic-resultant}

The four-cycle estimate has a counterpart for three labels on each
side. Here a polynomial Gram identity replaces the character table.
The bound applies to the complete bipartite graph $K_{3,3}$; deleting one of its
edges does not preserve the argument.

\begin{theorem}[Squarefree cubic resultant]\label{thm:cubic-resultant}
Let $p$ be an odd prime and let $\mathcal S_d$ be the monic squarefree
polynomials of degree $d$ over $\Fp$. On counting-measure spaces, set
\[
 B_d(f,g)=\chi(\operatorname{Res}(g,f)),
 \qquad f\in\mathcal S_d,\quad g\in\mathcal S_3.
\]
Then $B_3^*B_3\preceq p^3I$. Consequently the kernel
\[
 \prod_{i,j=1}^3\chi(y_j-x_i)
\]
on uniform $\Fp^3$, compressed on both sides to internally distinct
triples, has norm at most $6p^{-3/2}$.
\end{theorem}

\begin{proof}
We first record the polynomial character identity needed for the Gram.
For a monic squarefree polynomial $D$ of degree $2m$, put
$\chi_D(f)=\chi(\operatorname{Res}(D,f))$ and
$S_j(D)=\sum_{f\ {\rm monic},\,\deg f=j}\chi_D(f)$.
On the finite \'etale algebra $\Fp[X]/D$, use the additive character
$\psi_D(z)=\ee{\operatorname{res}(z/D)/p}$, where the residue is the
coefficient of $X^{-1}$ at infinity. In the monomial basis
$X^{2m-1},\ldots,X,1$, ordered by decreasing degree, the matrix of
$(u,v)\mapsto\operatorname{res}(uv/D)$ has antidiagonal entries $1$
and zeros below that antidiagonal, so its determinant is
$(-1)^{m(2m-1)}$. Factoring the algebra into fields and evaluating
the quadratic Gauss sums gives
\[
 \sum_z\chi_D(z)\psi_D(bz)=p^m\chi_D(b).
\]
This holds also for nonunits $b$: a field factor then has a zero
multiplicative-character sum. The annihilator of the polynomials
of degree less than $k$ consists of those of degree less than
$2m-k$. Summing the displayed Fourier identity over monic degree-$k$
polynomials, and then over the nonzero leading coefficient, yields
\[
 S_k(D)=p^{k-m}\left((p-1)\sum_{j=0}^{2m-k-2}S_j(D)
                         -S_{2m-k-1}(D)\right),\qquad k<2m.
\]
For $j\ge2m$, $S_j(D)=0$ by periodicity modulo $D$: the sum over
a complete residue system modulo $D$ vanishes.
In particular,
\[
\begin{array}{c|l}
 \deg D&\text{identities}\\ \hline
 6&S_3=(p-1)(1+S_1)-S_2,\\
 4&S_2=p-1-S_1,\quad S_3=-p,\\
 2&S_1=-1,\quad S_2=S_3=0.
\end{array}
\]

Write $G_d=B_d^*B_d$, let $r_g$ be the number of rational roots of
$g$, put $R=\diag(r_g)$ and $L=G_1+J$, and introduce the features
\[
 C(a,g)=\one[g(a)=0],\qquad
 V(a,t;g)=\one[g(a)=0]\chi\bigl((g/(X-a))(t)\bigr).
\]
Here $a,t\in\Fp$ and $g\in\mathcal S_3$; in $V$, the quotient is
understood only when $a$ is a root of $g$, and $V(a,t;g)=0$ otherwise.
Set $T=C^*C+V^*V$ and $E=\sum_aV(a,a)^*V(a,a)$.
The exact identity is
\begin{equation}\label{eq:cubic-gram}
 G_3=p^3I+RL+LR-L-G_2-(p+1)T+E.
\end{equation}
Here is an entrywise verification, including nonsplit cubics.
Fix $g,h$, put $U=\gcd(g,h)$, $D=gh/U^2$, and let $u$ be the number of
$a\in\Fp$ such that $g(a)h(a)=0$. With
$M_j=\sum_{\deg f=j,\ f\ {\rm monic}}\chi_g(f)\chi_h(f)$, the
nonsquarefree quadratics and cubics are, respectively,
$(X-a)^2$ and $(X-a)^2(X-b)$. Thus
\[
 (G_1)_{g,h}=M_1,\quad (G_2)_{g,h}=M_2-(p-u),\quad
 (G_3)_{g,h}=M_3-(p-u)M_1.
\]
For $g\ne h$, the generating series of the $M_j$ is that of the
$S_j(D)$ multiplied by $\prod_{P\mid U}(1-\chi_D(P)z^{\deg P})$.
If $U=1$, the degree-six identity gives zero correction in
\eqref{eq:cubic-gram}. If $U=X-a$, put
$\eta=\chi_D(X-a)$ and $\ell=1+S_1(D)$. Then
\[
 M_1=\ell-1-\eta,\quad M_2=p-(1+\eta)\ell+\eta,\quad
 M_3=-p(1+\eta)+\eta\ell,
\]
and $T_{g,h}=\ell$, $E_{g,h}=\eta$, giving the correction
$-(p+1)\ell+\eta$.
If $U=(X-a)(X-b)$, write
$e_1=\chi_D(X-a)+\chi_D(X-b)$,
$e_2=\chi_D(X-a)\chi_D(X-b)$. Then
$(M_1,M_2,M_3)=(-1-e_1,e_1+e_2,-e_2)$ and
$(T_{g,h},E_{g,h})=(-e_1,e_1)$, giving $(p+2)e_1$.
If $U$ is an irreducible quadratic, the same calculation with
$\eta=\chi_D(U)$ gives $(M_1,M_2,M_3)=(-1,-\eta,\eta)$ and
$T_{g,h}=E_{g,h}=0$.
On the diagonal, listing the degrees of the irreducible factors
of $g$, the three factorization types give
\[
\begin{array}{c|c|c|c|c|c}
 \text{type}&r_g&M_1&M_2&M_3&T_{g,g}\\ \hline
 111&3&p-3&p^2-3p+3&p^3-3p^2+3p-1&3(p-1)\\
 12&1&p-1&p^2-p-1&p^3-p^2-p+1&p+1\\
 3&0&p&p^2&p^3-1&0
\end{array}
\]
and $E_{g,g}=r_g$. Substitution proves \eqref{eq:cubic-gram}
in every case.

Finally $0\preceq E\preceq T$. The synthesis identities
\[
 \sum_a C(a,g)=r_g,\qquad
 \sum_a\chi(t-a)V(a,t;g)=r_g\chi(g(t))
\]
and Cauchy--Schwarz imply
$C^*C\succeq p^{-1}RJR$ and
$V^*V\succeq(p-1)^{-1}RG_1R$, hence $pT\succeq RLR$.
Since
$RLR-RL-LR+L=(R-I)L(R-I)\succeq0$,
\eqref{eq:cubic-gram} gives $G_3\preceq p^3I$.
Restriction to split cubics is a compression. Each split cubic has
six ordered root triples; lifting both variables therefore multiplies
the counting norm by $6$. Dividing by $p^3$ for uniform-cube measure
proves the last assertion.
\end{proof}

As with the four-cycle bound, separate input and output restrictions
are contractions. A jointly imposed missing-edge mask is a different
operator and needs its own proof. The complete nine-edge case is used
in the finite upper comparisons of \cref{prop:upper-soft-nine}.
Other three-by-three subgraphs there retain matching, path or
four-cycle witnesses.


\subsection{A two-variable character sum}\label{subsec:dl:H2}

The second estimate concerns two labels whose characters are coupled through
one factor $y-z$. For fixed $z$ the inner sum over $y$ is, up to an explicit
correction, minus the trace of Frobenius at $z$ on a sheaf of rank
$2\lfloor\deg f/2\rfloor$ attached to the family of curves $y^2=f(x)(x-z)$
(Katz--Sarnak). Twisting by the Kummer sheaf of $g$ keeps this sheaf
geometrically irreducible of rank at least two. So the $H^2_c$ coinvariants vanish (\cref{lem:tame-trace-bound} below), i.e.\ no term of order $p^{3/2}$ occurs; the compactly supported cohomology lives in degree one, its
dimension is the tame Euler characteristic, and Deligne's theorem bounds its
Frobenius eigenvalues by $p$. This is square-root cancellation in both
variables.

For an integer $m\ge1$ put $\gamma(m)=m$ if $m$ is even and $\gamma(m)=m-1$ if
$m$ is odd ($\gamma(m)=2\lfloor m/2\rfloor$; for $z$ distinct from the roots of $f$, this is twice the genus of the smooth projective model of $y^2=f(x)(x-z)$ with $\deg f=m$), and for integers $m\ge2$, $n\ge1$, $s\ge0$
\begin{equation}\label{eq:dl:B}
  \mathcal B(m,n,s)=\gamma(m)(m+n-s-1)+[m\text{ odd}]\big((n-1)r+mr^2\big)
  +m\big((m-2)r+r^2\big).
\end{equation}

\begin{theorem}[Two-variable character sum]\label{thm:dl:H2}
Let $p$ be an odd prime, $r=p^{-1/2}$, $m\ge2$ and $n\ge1$ integers,
$a_1,\dots,a_m\in\Fp$ distinct, $b_1,\dots,b_n\in\Fp$ distinct,
$c,c'\in\Fp^\times$,
\[
  f(X)=c\prod_{i=1}^m(X-a_i),\qquad g(X)=c'\prod_{k=1}^n(X-b_k),\qquad
  s=\big|\{a_1,\dots,a_m\}\cap\{b_1,\dots,b_n\}\big| .
\]
Then $\Sigma=\sum_{y,z\in\Fp}\chi\big(f(y)\,g(z)\,(y-z)\big)$ satisfies
$|\Sigma|\le p\,\mathcal B(m,n,s)$, and if $n\ge2$ also
$|\Sigma|\le p\,\mathcal B(n,m,s)$.
\end{theorem}

\noindent The main term $\gamma(m)(m+n-s-1)$ is the dimension of a cohomology
group of weight at most $2$; all other terms are the displayed
multiples of $r$ and $r^2$ in \eqref{eq:dl:B}, with coefficients
depending only on $m$ and $n$. (Only split
polynomials occur in this paper. For even $m$ the same proof works for
squarefree $f,g$ that do not split, with Weil's bound for general squarefree
polynomials in place of \cref{thm:weil}; for odd $m$ the construction in
\cref{lem:dl:sheaf} uses a root of $f$ in $\Fp$.) The proof rests on the
following lemma, in which $U_f=\mathbb A^1_{\Fp}\setminus\{a_1,\dots,a_m\}$ and
traces and Frobenius eigenvalues are regarded as complex numbers through the
isomorphism of $\overline{\mathbb Q}_\ell$ with $\C$ fixed in
\cref{subsec:external:sheaves}. Its parts (v) and (vi) are used only in
\cref{subsec:moving-root}.

\begin{lemma}[The sheaf of the family $y^2=f(x)(x-z)$]\label{lem:dl:sheaf}
Let $f=\prod_{i=1}^m(X-a_i)$ with $m\ge2$ and distinct $a_i\in\Fp$ \textup(so
$f$ is monic\textup), and put
\[
  \tau_f(z)=\frac1p\Big(\sum_{y\in\Fp}\chi\big(f(y)(y-z)\big)+[m\text{ odd}]\Big)
  \qquad(z\in\Fp).
\]
If $m$ is odd, fix a root $a$ of $f$ \textup(any root may be
chosen\textup). The construction below, applied to this fixed
$a$, gives a lisse $\overline{\mathbb Q}_\ell$-sheaf
$\mathcal F'$ on $U_f$ such that:
\begin{enumerate}[label=\textup{(\roman*)}]
\item $\mathcal F'$ has rank $\gamma(m)$ and is pure of weight $1$;
\item $\mathcal F'$ is geometrically irreducible;
\item $\mathcal F'$ is tame at every point of
  $\mathbb P^1\setminus U_f=\{a_1,\dots,a_m,\infty\}$;
\item $p\,\tau_f(z)=-\mathrm{Tr}(\mathrm{Frob}_z\mid\mathcal F')$ for every
  $z\in U_f(\Fp)$;
\item there is a perfect alternating pairing
  $\mathcal F'\otimes\mathcal F'\to\overline{\mathbb Q}_\ell(-1)$ of lisse sheaves
  on $U_f$;
\item for every root $a_i$ of $f$ if $m$ is even, and for every root $a_i\ne a$
  if $m$ is odd, the inertia group at $a_i$ acts on $\mathcal F'$ through
  unipotent automorphisms, not all trivial.
\end{enumerate}
\end{lemma}

\begin{proof}
\emph{$m$ even.} Let $\mathcal F'=\mathcal F_f$ be the sheaf of
\cref{thm:prelim:ks} with $k=\Fp$ and $h=f$, a polynomial of degree
$m=2\cdot\frac m2$, $\frac m2\ge1$, with $m$ distinct roots; then $U_h=U_f$.
Parts (i), (ii), (iii) and (v) are \cref{thm:prelim:ks}(i), (iii), (ii) and
(v), and \cref{thm:prelim:ks}(iv) with $E=\Fp$, $t=z$, gives
$-\mathrm{Tr}(\mathrm{Frob}_z\mid\mathcal F_f)=\sum_x\chi((x-z)f(x))=p\,\tau_f(z)$.
For (vi), $\mathcal F_f$ is tame at $a_i$ (\cref{thm:prelim:ks}(ii)), so its
local monodromy there factors through the pro-cyclic tame quotient
(\ref{it:prelim:E3}), and the inertia image is the closure of the group
generated by the element $u$ of \cref{thm:prelim:ks}(vi). As $(u-1)^2=0$, the
powers $u^k=1+k(u-1)$ are unipotent, hence so is every element
$1+c(u-1)$ of the closure, and $u\ne1$.

\emph{$m$ odd.} Then $m\ge3$. Put $f_1=f/(X-a)$, so $\deg f_1=m-1\ge2$ is even
and $f_1(a)\ne0$, and
\[
  f_1^\sharp(X)=X^{m-1}f_1\big(a+X^{-1}\big)=\prod_{a_i\ne a}\big(1+(a-a_i)X\big).
\]
$f_1^\sharp$ has degree $m-1$ (leading coefficient $f_1(a)\ne0$), the $m-1$
distinct roots $1/(a_i-a)\in\Fp^\times$ ($a_i\ne a$), and $f_1^\sharp(0)=1$.
Let $\varphi_a:\mathbb A^1\setminus\{a\}\to\mathbb A^1\setminus\{0\}$,
$\varphi_a(z)=1/(z-a)$, an isomorphism of $\Fp$-schemes that extends to the
automorphism $z\mapsto1/(z-a)$ of $\mathbb P^1$ (with $a\mapsto\infty$,
$\infty\mapsto0$). It maps $U_f$ onto
$W=\mathbb A^1\setminus(\{0\}\cup Z(f_1^\sharp))$, since
$\varphi_a(a_i)=1/(a_i-a)$ for $a_i\ne a$. Let $\mathcal F_{f_1^\sharp}$ be the
sheaf of \cref{thm:prelim:ks} with $k=\Fp$, $h=f_1^\sharp$ (degree
$m-1=2\cdot\frac{m-1}2$, $\frac{m-1}2\ge1$, distinct roots), lisse on
$U_{f_1^\sharp}=\mathbb A^1\setminus Z(f_1^\sharp)\supseteq W$, and let
$\mathcal L=\mathcal L_{\chi(a-z)}$ be the Kummer sheaf \ref{it:prelim:E4} of the
polynomial $a-z$, lisse on $\mathbb A^1\setminus\{a\}\supseteq U_f$. Put
\[
  \mathcal F'=\varphi_a^*\big(\mathcal F_{f_1^\sharp}|_W\big)\otimes\mathcal L|_{U_f}.
\]
\emph{Transport of (i)--(iii) and (v).} Restriction to a dense open,
pull-back by an isomorphism, and tensoring with a rank-one sheaf preserve
geometric irreducibility (\ref{it:prelim:E1}). Thus
\cref{thm:prelim:ks}(i),(iii) and \ref{it:prelim:E2}, \ref{it:prelim:E4}
give rank $m-1=\gamma(m)$, purity of weight $1$ and geometric
irreducibility for $\mathcal F'$. For tameness, the automorphism
$\varphi_a$ sends the roots other than $a$ to $Z(f_1^\sharp)$, sends $a$
to infinity, and sends infinity to $0$. The original sheaf is tame at
$Z(f_1^\sharp)$ and infinity and lisse at $0$; hence its restriction to $W$
and pull-back are tame at every missing point. The Kummer sheaf is tame
there too, so their tensor product is tame (\ref{it:prelim:E3}). Finally,
$\mathcal L\otimes\mathcal L$ is the constant sheaf, since its monodromy
character has order dividing two (\ref{it:prelim:E4}). Pulling back the
pairing of \cref{thm:prelim:ks}(v) and using this trivialization gives the
perfect alternating pairing for $\mathcal F'$, with the same target
$\overline{\mathbb Q}_\ell(-1)$. This proves (i)--(iii) and (v).

(iv) Let $z\in U_f(\Fp)$ and $w=\varphi_a(z)=1/(z-a)$, so $-w=1/(a-z)$. By
\ref{it:prelim:E1}, \ref{it:prelim:E4} and \cref{thm:prelim:ks}(iv) with $E=\Fp$,
\[
  -\mathrm{Tr}(\mathrm{Frob}_z\mid\mathcal F')
  =-\chi(a-z)\,\mathrm{Tr}(\mathrm{Frob}_w\mid\mathcal F_{f_1^\sharp})
  =\sum_{x'\in\Fp}\chi\big((a-z)(x'-w)f_1^\sharp(x')\big).
\]
On the other hand, the term $y=a$ of $\sum_y\chi(f(y)(y-z))$ vanishes
($f(a)=0$), and $y\mapsto x'=1/(y-a)$ is a bijection of $\Fp\setminus\{a\}$ onto
$\Fp^\times$ with $y=a+1/x'$. For such $y$,
$f(y)=x'^{-1}f_1(a+1/x')=x'^{-m}f_1^\sharp(x')$ and
$y-z=x'^{-1}\big((a-z)x'+1\big)=x'^{-1}(a-z)(x'-w)$, so
$f(y)(y-z)=x'^{-(m+1)}(a-z)(x'-w)f_1^\sharp(x')$ with $m+1$ even. Hence
\begin{align*}
  \sum_y\chi\big(f(y)(y-z)\big)
  &=\sum_{x'\ne0}\chi\big((a-z)(x'-w)f_1^\sharp(x')\big)\\
  &=-\mathrm{Tr}(\mathrm{Frob}_z\mid\mathcal F')-\chi\big((a-z)(-w)f_1^\sharp(0)\big)
  =-\mathrm{Tr}(\mathrm{Frob}_z\mid\mathcal F')-1,
\end{align*}
since $(a-z)(-w)=1$ and $f_1^\sharp(0)=1$; that is,
$p\,\tau_f(z)=-\mathrm{Tr}(\mathrm{Frob}_z\mid\mathcal F')$.
(vi) Let $a_i\ne a$ be a root of $f$. The automorphism of $\mathbb P^1$
extending $\varphi_a$ maps $a_i$ to the root $1/(a_i-a)$ of $f_1^\sharp$ and,
by \ref{it:prelim:E3}, the inertia groups at $a_i$ onto those at $1/(a_i-a)$;
so the inertia at $a_i$ acts on $\varphi_a^*(\mathcal F_{f_1^\sharp}|_W)$ as the
inertia at $1/(a_i-a)$ acts on $\mathcal F_{f_1^\sharp}|_W$, that is, as it acts
on $\mathcal F_{f_1^\sharp}$ (restriction to the dense open $W$ keeps the
inertia groups and the local monodromy at this point, \ref{it:prelim:E3}).
This is as stated in (vi), by the argument of
the even case applied to $\mathcal F_{f_1^\sharp}$ (\cref{thm:prelim:ks}(ii),
(vi)). The sheaf $\mathcal L$ is lisse at $a_i$, so the
inertia at $a_i$ acts trivially on it (\ref{it:prelim:E3}), and the tensor
product has the same local monodromy at $a_i$ as
$\varphi_a^*\mathcal F_{f_1^\sharp}$.
\end{proof}

For odd $m$ we have $m\ge3$, so the root $a$ can be chosen different from any
given root of $f$. Thus every root of $f$ is, for some allowed
choice of the root $a$ in the construction of $\mathcal F'$, a
point at which the inertia acts through nontrivial unipotent
automorphisms; this is how part (vi) is used in
\cref{lem:moving-root-moments}.

The cohomological step used below is the same for one trace and for a product
of two traces. Once the geometric coinvariants vanish, only the rank, the
weight and the number of punctures enter the bound.

\begin{lemma}[Tame trace bound]\label{lem:tame-trace-bound}
Let $B\subseteq\Fp$ have $m\ge1$ elements, let
$U=\mathbb A^1_{\Fp}\setminus B$, and let $\mathcal G$ be a lisse sheaf on $U$
of rank $R$, mixed of weights at most $w$, and tame at every point of
$B\cup\{\infty\}$. Suppose its generic stalk $V$ has zero coinvariants under
$\pi_1^{\rm geom}(U)$. Then
\[
 \dim H^1_c(U\otimes\bar\F_p,\mathcal G)=R(m-1),\qquad
 \left|\sum_{x\in U(\Fp)}\mathrm{Tr}(\mathrm{Frob}_x\mid\mathcal G)\right|
 \le R(m-1)p^{(w+1)/2}.
\]
In particular the coinvariant hypothesis holds if $\mathcal G$ is
geometrically irreducible of rank at least two.
\end{lemma}

\begin{proof}
By \cref{thm:prelim:curves}, compactly supported cohomology vanishes outside
degrees one and two, and $H^2_c\cong V_{\pi_1^{\rm geom}(U)}(-1)=0$.
Tameness at all $m+1$ missing points makes every Swan conductor zero, so
\cref{thm:prelim:gos} gives
$-\dim H^1_c=R(2-(m+1))$. By \cref{thm:prelim:lefschetz}, the trace sum is
$-\mathrm{Tr}(\mathrm{Frob}_p\mid H^1_c)$, and every eigenvalue in this trace
has modulus at most $p^{(w+1)/2}$ by \cref{thm:prelim:weights}.
Finally, the span of the vectors $gv-v$ is invariant. For an irreducible
representation it is zero or all of $V$; the first possibility makes the
action trivial, which is incompatible with irreducibility in rank at least
two. Thus the coinvariants vanish in the stated special case.
\end{proof}

\begin{proof}[Proof of \cref{thm:dl:H2}]
\emph{Step 0 (normalization).} Replacing $f,g$ by $f/c$, $g/c'$ multiplies
$\Sigma$ by $\chi(cc')=\pm1$, and $\mathcal B$ does not depend on $c,c'$; so we
may assume $f$ and $g$ monic. Write $Z(f)=\{a_i\}$, $Z(g)=\{b_k\}$, so
$|Z(f)\cup Z(g)|=m+n-s$.

\emph{Step 1 (the three parts).}
$\Sigma=\sum_{z\in\Fp}\chi(g(z))\sum_y\chi\big(f(y)(y-z)\big)$. For
$z=a_j\in Z(f)$, $f(y)(y-a_j)=(y-a_j)^2f_j(y)$ with
$f_j=\prod_{i\ne j}(X-a_i)$, $\deg f_j=m-1\ge1$, so
$\chi(f(y)(y-a_j))=\chi(f_j(y))$ for $y\ne a_j$ and $0$ for $y=a_j$; by
\cref{thm:weil},
$|\sum_y\chi(f(y)(y-a_j))|\le|\sum_y\chi(f_j(y))|+1\le(m-2)\sqrt p+1$. For
$z\notin Z(f)$, \cref{lem:dl:sheaf}(iv) gives
$\sum_y\chi(f(y)(y-z))=p\,\tau_f(z)-[m\text{ odd}]
=-\mathrm{Tr}(\mathrm{Frob}_z\mid\mathcal F')-[m\text{ odd}]$. Hence
\begin{multline*}
  \Sigma=-\sum_{z\in U_f(\Fp)}\chi(g(z))\,\mathrm{Tr}(\mathrm{Frob}_z\mid\mathcal F')
  \ -\ [m\text{ odd}]\sum_{z\in U_f(\Fp)}\chi(g(z))\\
  +\ \sum_{z\in Z(f)}\chi(g(z))\sum_y\chi\big(f(y)(y-z)\big).
\end{multline*}
The last part is at most $m\big((m-2)\sqrt p+1\big)$ in modulus. The middle sum
differs from $\sum_{z\in\Fp}\chi(g(z))$ by at most $m$ terms of modulus at most
$1$, and $|\sum_z\chi(g(z))|\le(n-1)\sqrt p$ by \cref{thm:weil}; so the middle
part is at most $(n-1)\sqrt p+m$ (and occurs only for odd $m$).

\emph{Step 2 (the sheaf $\mathcal G$).} Let
$U'=U_f\setminus Z(g)=\mathbb A^1\setminus(Z(f)\cup Z(g))$,
$\mathcal L_{\chi(g)}$ the Kummer sheaf of $g$ of \ref{it:prelim:E4}, lisse on
$\mathbb A^1\setminus Z(g)\supseteq U'$, and
$\mathcal G=\mathcal L_{\chi(g)}|_{U'}\otimes\mathcal F'|_{U'}$. As $\chi(g(z))=0$
for $z\in Z(g)$, the first part is
$-\sum_{z\in U'(\Fp)}\mathrm{Tr}(\mathrm{Frob}_z\mid\mathcal G)$ by
\ref{it:prelim:E1} and \ref{it:prelim:E4}. $\mathcal G$ is lisse of rank
$\gamma(m)\ge2$ and pure of weight $1$ (\cref{lem:dl:sheaf}(i),
\ref{it:prelim:E2}); it is geometrically irreducible (\cref{lem:dl:sheaf}(ii):
the restriction to the dense open $U'$ and the tensor product with a rank-one
sheaf preserve this, \ref{it:prelim:E1}); and it is tame at every point of
$\mathbb P^1\setminus U'=Z(f)\cup Z(g)\cup\{\infty\}$ by \ref{it:prelim:E3}
(restriction to a dense open keeps tameness, a sheaf is tame at a point where it
is lisse, and tensor products of tame sheaves are tame): $\mathcal F'$ is tame at
$Z(f)\cup\{\infty\}$ (\cref{lem:dl:sheaf}(iii)) and lisse at the points of
$Z(g)\setminus Z(f)\subseteq U_f$, and $\mathcal L_{\chi(g)}$ is tame at the
points of $Z(g)\cup\{\infty\}$ (\ref{it:prelim:E4}) and lisse at those of
$Z(f)\setminus Z(g)$.

\emph{Step 3 (cohomology).} The sheaf of Step~2 satisfies
\cref{lem:tame-trace-bound}: its rank is $\gamma(m)\ge2$, its weight is $1$,
it is geometrically irreducible, and it is tame at all
$m+n-s+1$ missing points, including infinity. Therefore
\[
 \left|\sum_{z\in U'(\Fp)}\mathrm{Tr}(\mathrm{Frob}_z\mid\mathcal G)\right|
 \le\gamma(m)(m+n-s-1)p.
\]
The trace sum is $-p\sum_{z\in U'(\Fp)}\chi(g(z))\tau_f(z)$, a rational
integer by \cref{lem:dl:sheaf}(iv), so the bound is independent of the
identification of $\overline{\mathbb Q}_\ell$ with $\C$. It bounds the first
part of Step~1.

\emph{Step 4 (conclusion).} Adding the three parts,
\[
  |\Sigma|\le\gamma(m)(m+n-s-1)p+[m\text{ odd}]\big((n-1)\sqrt p+m\big)
  +m\big((m-2)\sqrt p+1\big)=p\,\mathcal B(m,n,s).
\]
Finally, $\Sigma=\chi(-1)\sum_{z,y}\chi\big(g(z)f(y)(z-y)\big)$ is, up to sign,
the sum of the theorem for the pair $(g,f)$; if $n\ge2$ the bound just proved
applies to it and gives $|\Sigma|\le p\,\mathcal B(n,m,s)$.
\end{proof}

\begin{remark}[Sharpness and the coefficients of $\mathcal B$]\label{rem:dl:H2-sharp}
(i) The main term is nearly attained (numerically): for $f$
proportional to $g$ of degree two
($m=n=s=2$) the group $H^1_c$ has dimension $2$, and exact computations for
primes up to $9\cdot10^4$ give $|\Sigma|/(p\mathcal B)$ up to $0.998$
(numerical, not used). The consistency test of \cref{thm:dl:H2} against exact
double sums is described in \cref{rem:cert:deligne-tests}; it is not an input
of any proof.
(ii) As a polynomial in $r$, $\mathcal B(m,n,0)=C_0+C_1r+C_2r^2$ with
$(C_0,C_1,C_2)=\big(\gamma(m)(m+n-1),\ m(m-2)+[m\text{ odd}](n-1),\
m(1+[m\text{ odd}])\big)$; these are the coefficient triples with which the
completion variance (\cref{lem:push:localvariance}) uses the theorem.
(iii) $\mathcal B(m,n,s)$ is nonincreasing in $s$ and nondecreasing in $n$.
Since $\mathcal B(2,4,0)=10+2r^2$, $\mathcal B(3,4,0)=12+6r+6r^2$ and
$\mathcal B(4,4,0)=28+8r+4r^2$, and $r\le1/\sqrt3$, we get
$\mathcal B(m,n,s)\le\mathcal B(4,4,0)<40$ for $2\le m\le4$, $1\le n\le4$,
$s\ge0$ and every odd prime; the product rows use this.
\end{remark}

\subsection{Trace functions with one free label}\label{subsec:moving-root}

At a root cut of the upper range the star part of the cut kernel is an average
over the label $z$ of an integrated centre of a product of characters
$\chi(z-x_j)$ over the leaves $j$ of the star. Fixing all half labels except
the label $x_i$ of one leaf turns such a star into an affine function of a
single \emph{moving-root trace} $\tau_S(x_i)$, whose fixed roots $S$ are the
labels of the other leaves (\cref{lem:star-trace}). The second moments of
combinations of these traces against the conditional law of $x_i$ (a
uniform law on an allowed set, whose density is a product of characters with
the neighbouring labels as roots) are governed by the almost-orthogonality of
traces with distinct root sets against character weights
(\cref{lem:moving-root-moments}). This is a statement about sheaves: distinct
root sets give non-isomorphic sheaves, even after a quadratic twist, because
their local monodromy at a root in the symmetric difference differs. First moments satisfy \eqref{eq:advance:rank-first}, with its factor $r^2$ in place of the pointwise factor $r$ in $|\tau_S|\le sr$ (\cref{lem:moving-root-trace}(b)), and they make pivot groups almost
orthogonal without any independence assumption
(\cref{prop:advance:ordered-star}). For a graph law the same conditioning,
with the conditional means bounded through \cref{thm:dl:H2} and the diagonal
terms through densities, bounds star marginals with arbitrary real
coefficients (\cref{lem:wt:stars}). The subsection ends with the exact mean of
the complete star (\cref{lem:star-mean}).

\subsubsection{Trace functions and stars}\label{subsec:advance:conditional-star}

For a set $S\subseteq\Fp$ of $s=|S|\ge2$ distinct elements put
$f_S=\prod_{c\in S}(X-c)$ and
\begin{equation}\label{eq:moving-root-tau}
  \tau_S(x)=\frac1p\Big(\sum_{z\in\Fp}\chi\Big((z-x)\prod_{c\in S}(z-c)\Big)
  +\one_{|S|\ \rm odd}\Big)\qquad(x\in\Fp).
\end{equation}
Thus $\tau_S=\tau_{f_S}$ in the notation of \cref{lem:dl:sheaf}. We write
$\mathcal F_S$ for a sheaf on $\mathbb A^1\setminus S$ given by
\cref{lem:dl:sheaf} for $f=f_S$ (for odd $s$ with some choice of the root $a$),
and $h(s)=2\lfloor s/2\rfloor=\gamma(s)$ for its rank. The correction
$\one_{|S|\ \rm odd}$ accounts for the points at infinity: for
$x\notin S$ and odd $s$ the polynomial $(z-x)f_S(z)$ is squarefree of even
degree with leading coefficient one, so the smooth projective curve
$y^2=(z-x)f_S(z)$ has two rational points at infinity, which the affine sum
misses. With the correction, $p\,\tau_S(x)$ is, for every $s$ and $x\notin S$,
minus the trace of Frobenius on the first cohomology of that curve.

\begin{lemma}[Moving-root traces]\label{lem:moving-root-trace}
Let $S\subseteq\Fp$ with $|S|=s\ge2$.
\begin{enumerate}[label=\textup{(\alph*)}]
\item $\tau_S$ is real, and $\tau_S(x)=-p^{-1}\,\mathrm{Tr}(\mathrm{Frob}_x\mid\mathcal F_S)$
  for $x\notin S$.
\item $|\tau_S(x)|\le h(s)r$ for $x\notin S$, and
  $|\tau_S(x)|\le\big((s-2)\sqrt p+2\big)/p\le sr$ for $x\in S$. In particular
  $|\tau_S|\le sr$ on $\Fp$.
\end{enumerate}
\end{lemma}

\begin{proof}
(a) is \cref{lem:dl:sheaf}(iv). (b) For $x\notin S$, $\mathcal F_S$ has rank
$h(s)$ and is pure of weight $1$ (\cref{lem:dl:sheaf}(i)), so each Frobenius
eigenvalue at $x$ has modulus $\sqrt p$, and $|p\,\tau_S(x)|\le h(s)\sqrt p$.
For $x\in S$, $(z-x)f_S(z)=(z-x)^2f_{S\setminus\{x\}}(z)$; the term $z=x$
vanishes, and for $z\ne x$ the summand is $\chi(f_{S\setminus\{x\}}(z))$. Hence
$p\,\tau_S(x)=\sum_z\chi(f_{S\setminus\{x\}}(z))-\chi(f_{S\setminus\{x\}}(x))
+\one_{s\ \rm odd}$. The first sum has modulus at most $(s-2)\sqrt p$ by
\cref{thm:weil} (a product of $s-1\ge1$ distinct linear factors), and the two
other terms have modulus at most one; finally $2\le2\sqrt p$.
\end{proof}

\begin{lemma}[Star-trace dictionary]\label{lem:star-trace}
Let $x:U\to\Fp$ be injective, $T\subseteq U$ nonempty, and
$w_T(x)=\E_z\prod_{j\in T}\chi(z-x_j)$. Then $w_T(x)=0$ if $|T|=1$,
$w_T(x)=-1/p$ if $|T|=2$, and for $|T|\ge3$ and every $i\in T$,
\[
  w_T(x)=\tau_{x(T\setminus\{i\})}(x_i)-\frac{\one_{|T|\ \rm even}}p,\qquad
  x(T\setminus\{i\})=\{x_j:j\in T\setminus\{i\}\}.
\]
For signs $\varsigma_j\in\{\pm1\}$,
$\E_z\prod_{j\in T}\chi(\varsigma_j(z-x_j))=\big(\prod_{j\in T}\chi(\varsigma_j)\big)w_T(x)$.
\end{lemma}

\begin{proof}
For $|T|=1$, $\sum_z\chi(z-x_j)=0$. For $|T|=2$, $\sum_z\chi((z-x_j)(z-x_{j'}))=-1$
by \eqref{eq:two-root}, as $x_j\ne x_{j'}$. For $|T|\ge3$ and $i\in T$,
$\prod_{j\in T}(z-x_j)=(z-x_i)\prod_{c\in S}(z-c)$ with
$S=x(T\setminus\{i\})$, a set of $|T|-1\ge2$ distinct elements, and
$|S|$ is odd iff $|T|$ is even; so $w_T(x)=\tau_S(x_i)-\one_{|T|\ \rm even}/p$
by \eqref{eq:moving-root-tau}. The signs factor out since $\chi$ is
multiplicative.
\end{proof}

\paragraph{The star function.} Consider spectral data (\cref{def:spec:data})
with half vertex set $U^+$, half graph $G^+$, cross signs $\varsigma_{jk}$ and
$N(k)=\{j\in U^+:\varsigma_{jk}\ne0\}$. The column stars of the cut kernel
centred at $k$ have kernels $\prod_{j\in T}\chi(\varsigma_{jk}(x'_k-x_j))$,
$\emptyset\ne T\subseteq N(k)$; averaging the centre label $z$ uniformly on $\Fp$ gives the \emph{integrated-centre star function}
\begin{equation}\label{eq:star-function}
  h_\star(x)=\sum_{k\in U^+}\Big(\E_z\prod_{j\in N(k)}\big(1+\chi(\varsigma_{jk}(z-x_j))\big)-1\Big)
  =\sum_{\emptyset\ne T\subseteq U^+}a_Tw_T(x),\qquad
  a_T=\sum_{k:\,T\subseteq N(k)}\ \prod_{j\in T}\chi(\varsigma_{jk}).
\end{equation}
We fix this sign convention: the signs $\chi(\varsigma_{jk})$ belong to the
centre $k$, and no factor $\chi(-1)^{|T|}$ appears. (Writing the characters as
$\chi(x_j-z)$ multiplies $a_T$ by $\chi(-1)^{|T|}$; the bounds below use $|a_T|$
for odd $|T|$, and for even $|T|$ the factor is one, so they do not depend on
the convention.) By \cref{lem:star-trace}, on injective labellings $h_\star$ is a
constant plus a linear combination of the traces
$\tau_{x(T\setminus\{i\})}(x_i)$, $|T|\ge3$, for any choice of $i\in T$. The
row stars map constants to constants
(\cref{lem:spec:stars}(b)), so they disappear under projection off constants.

\subsubsection{Means and mean squares of the traces $\tau_S$}\label{subsec:moving-root-moments}

For a finite set $A\subseteq\Fp$ put $\chi_A(x)=\prod_{a\in A}\chi(x-a)$, the
squarefree character factor with root set $A$ ($\chi_\emptyset=1$).

\begin{lemma}[Moving-root moments]\label{lem:moving-root-moments}
Let $p$ be an odd prime, $r=p^{-1/2}$, $h(s)=2\lfloor s/2\rfloor$, and let
$S,T,A\subseteq\Fp$ with $|S|=s\ge2$, $|T|=t\ge2$ and $|A|=u\ge0$; the three
sets may overlap arbitrarily. Then
\begin{align}
 \Big|\E_x\chi_A\tau_S\tau_T-p^{-1}\one_{S=T,\ A=\emptyset}\Big|
 &\le r^3\Big[(s+t+u-1)h(s)h(t)+1+r(s+t+u)st\Big],
 \label{eq:advance:rank-second}\\
 \big|\E_x\chi_A\tau_S\big|
 &\le r^2\Big[(s+u-1)h(s)+r(s+u)s\Big].
 \label{eq:advance:rank-first}
\end{align}
\end{lemma}

\begin{proof}
Let $U=\mathbb A^1\setminus(S\cup T\cup A)$ and $\pi=\pi_1^{\rm geom}(U)$. We use
sheaves $\mathcal F_S$, $\mathcal F_T$ of \cref{lem:dl:sheaf} for $f_S$, $f_T$,
restricted to $U$, with $\mathcal F_T=\mathcal F_S$ if $T=S$; when $s$ or $t$ is
odd, the root $a$ of that construction (the \emph{inversion root}) is chosen in
Step~3 (\cref{lem:dl:sheaf} and \cref{lem:moving-root-trace} hold for every
choice). Let $\mathcal L$ be the Kummer sheaf of
$g_A=\prod_{a\in A}(X-a)$ restricted to $U$ (the constant sheaf if
$A=\emptyset$), with monodromy character $\lambda$; so
$\mathrm{Tr}(\mathrm{Frob}_x\mid\mathcal L)=\chi_A(x)$ for $x\in U(\Fp)$
(\ref{it:prelim:E4}). Put $\mathcal G=\mathcal F_S\otimes\mathcal F_T\otimes\mathcal L$,
and write $V_S$, $V_T$ for the generic stalks of $\mathcal F_S$, $\mathcal F_T$ as
representations $\rho_S$, $\rho_T$ of $\pi$.

\emph{Step 1 (traces and omitted points).} For every $x\in U(\Fp)$,
\cref{lem:moving-root-trace}(a) and \ref{it:prelim:E1} give
\[
  \chi_A(x)\tau_S(x)\tau_T(x)=p^{-2}\,\mathrm{Tr}(\mathrm{Frob}_x\mid\mathcal G).
\]
For $x\in A$, $\chi_A(x)=0$. The remaining points $x\in(S\cup T)\setminus A$ number at
most $s+t$, and there $|\chi_A\tau_S\tau_T|\le st\,r^2$ by
\cref{lem:moving-root-trace}(b). Hence
\begin{equation}\label{eq:moving-root-omitted}
  \Big|\E_x\chi_A\tau_S\tau_T-p^{-3}\sum_{x\in U(\Fp)}
  \mathrm{Tr}(\mathrm{Frob}_x\mid\mathcal G)\Big|
  \le p^{-1}(s+t)st\,r^2\le r^4(s+t+u)st .
\end{equation}

\emph{Step 2 (the top cohomology).} $\mathcal G$ is lisse on $U$ of rank
$h(s)h(t)$, pure of weight $2$ (\ref{it:prelim:E2}; $\mathcal L$ has weight $0$),
and tame at every point of $\mathbb P^1\setminus U$ (\cref{lem:dl:sheaf}(iii),
\ref{it:prelim:E3}, \ref{it:prelim:E4}). By
\cref{thm:prelim:lefschetz,thm:prelim:curves},
\[
  \sum_{x\in U(\Fp)}\mathrm{Tr}(\mathrm{Frob}_x\mid\mathcal G)
  =\mathrm{Tr}(\mathrm{Frob}_p\mid H^2_c)-\mathrm{Tr}(\mathrm{Frob}_p\mid H^1_c),
  \qquad H^i_c=H^i_c(U\otimes\bar\F_p,\mathcal G),
\]
and $H^2_c\cong(V_S\otimes V_T\otimes\lambda)_\pi(-1)$. The dual of the space of
coinvariants is the space of $\pi$-invariant linear forms on
$V_S\otimes V_T\otimes\lambda$, that is the $\pi$-invariants of
$V_S^\vee\otimes V_T^\vee\otimes\lambda^{-1}$. The pairing of
\cref{lem:dl:sheaf}(v) identifies $V_S$ with $V_S^\vee$ as $\pi$-representations
($\pi$ acts trivially on $\overline{\mathbb Q}_\ell(-1)$), and $\lambda^{-1}=\lambda$;
so this space is $(V_S\otimes V_T^\vee\otimes\lambda)^\pi
=\mathrm{Hom}_\pi(V_T,V_S\otimes\lambda)$. The representations $V_T$ and
$V_S\otimes\lambda$ are irreducible (\cref{lem:dl:sheaf}(ii); restriction to a
dense open and twisting by a rank-one sheaf preserve geometric irreducibility,
\ref{it:prelim:E1}). By Schur's lemma, $\dim H^2_c=1$ if
$V_S\cong V_T\otimes\lambda$ (equivalently $V_T\cong V_S\otimes\lambda$, as
$\lambda^2=1$), and $H^2_c=0$ otherwise.

\emph{Step 3 (local monodromy: $H^2_c=0$ unless $S=T$ and
$A=\emptyset$).} For a finite point $c\in\mathbb P^1\setminus U$ let
$I_c\subseteq\pi$ be an inertia group at $c$ (\ref{it:prelim:E3}). A sheaf lisse
at $c$ has trivial inertia image there, and a rank-one sheaf has scalar inertia
image. Restriction to the dense open $U$ keeps the local monodromy at $c$
(\ref{it:prelim:E3}), so \cref{lem:dl:sheaf}(vi) for $\mathcal F_S$ on
$\mathbb A^1\setminus S$ holds for its restriction to $U$. For every $c\in A$
the character $\lambda$ is nontrivial on $I_c$: $c$ is a simple root of $g_A$, so
the double covering $w^2=g_A(z)$ is ramified at $c$ and the local monodromy of
$\mathcal L$ at $c$ is the character of order $2$ of $I_c$ (the case of odd
multiplicity in \ref{it:prelim:E4}, and \ref{it:prelim:E3} for the restriction
to $U$).

Suppose $S\ne T$. Choose $c_0\in S\triangle T$; exchanging $S$ and $T$ if
necessary, $c_0\in S\setminus T$. If $s$ is odd, construct $\mathcal F_S$ with an
inversion root different from $c_0$. By \cref{lem:dl:sheaf}(vi),
$\rho_S(I_{c_0})$ is not contained in the scalars (a unipotent scalar is the
identity), while $\mathcal F_T$ is lisse at $c_0$ and $\mathcal L$ has rank one,
so $\rho_T\otimes\lambda$ maps $I_{c_0}$ into the scalars. Hence
$V_S\not\cong V_T\otimes\lambda$.

Suppose $S=T$ and $A\ne\emptyset$, so that $\mathcal F_T=\mathcal F_S$. Choose
$c_0\in A$ and, if $c_0\in S$ and $s$ is odd, construct $\mathcal F_S$ with an
inversion root different from $c_0$. Take $g\in I_{c_0}$ with $\lambda(g)=-1$.
If $c_0\notin S$, then $\rho_S(g)=1$ ($\mathcal F_S$ is lisse at $c_0$), while
$(\rho_S\otimes\lambda)(g)=-1$. If $c_0\in S$, then $\rho_S(g)$ is unipotent by
\cref{lem:dl:sheaf}(vi), so all its eigenvalues are $1$, while all eigenvalues of
$(\rho_S\otimes\lambda)(g)=-\rho_S(g)$ are $-1$. In both cases $\rho_S(g)$ and
$(\rho_S\otimes\lambda)(g)$ are not conjugate, so $V_S\not\cong V_S\otimes\lambda$.

Thus $H^2_c=0$ unless $S=T$ and $A=\emptyset$. (In that case $\lambda=1$,
$V_T=V_S$ and $H^2_c$ is one-dimensional; the case is computed directly in
Step~4.)

\emph{Step 4 (the diagonal case).} Let $S=T$ and $A=\emptyset$. Write
$p\,\tau_S(x)=Y(x)+\one_{s\ \rm odd}$ with $Y(x)=\sum_z\chi(z-x)\chi(f_S(z))$.
Then $\sum_xY(x)=0$, since $\sum_x\chi(z-x)=0$. By the two-root identity
\eqref{eq:two-root}, $\sum_x\chi\big((z-x)(z'-x)\big)$ is $p-1$ if $z=z'$ and
$-1$ otherwise, and $\chi(f_S(z))^2=\one_{z\notin S}$; so
\begin{gather*}
  \sum_xY(x)^2=p(p-s)-\Big(\sum_z\chi(f_S(z))\Big)^2,\\
  \E_x\tau_S^2=\frac{\sum_xY(x)^2+p\one_{s\ \rm odd}}{p^3}
  =\frac1p-\frac{h(s)}{p^2}-\frac1{p^3}\Big(\sum_z\chi(f_S(z))\Big)^2,
\end{gather*}
as $s-\one_{s\ \rm odd}=h(s)$. By \cref{thm:weil}(a) the last sum has modulus at
most $(s-1)\sqrt p$, so
$|\E_x\tau_S^2-p^{-1}|\le\big(h(s)+(s-1)^2\big)r^4\le2s^3r^4$, which is the last
term of \eqref{eq:advance:rank-second} for $T=S$, $u=0$. Conceptually, the main
term $1/p$ is the one invariant of $\mathcal F_S\otimes\mathcal F_S$, the pairing
of \cref{lem:dl:sheaf}(v), on which Frobenius acts by $p$; the computation does
not use this. The sheaves are needed for distinct root sets and nontrivial
weights.

\emph{Step 5 (dimension and weights).} If $S\ne T$ or $A\ne\emptyset$,
Steps~2 and~3 show that the geometric coinvariants of $\mathcal G$ vanish.
Its rank is $h(s)h(t)$, its weight is $2$, and it is tame at the
$|S\cup T\cup A|+1$ missing points, including infinity. Thus
\cref{lem:tame-trace-bound} gives
\[
 \Big|\sum_{x\in U(\Fp)}\mathrm{Tr}(\mathrm{Frob}_x\mid\mathcal G)\Big|
 \le h(s)h(t)(|S\cup T\cup A|-1)p^{3/2}
 \le(s+t+u-1)h(s)h(t)p^{3/2}.
\]
Dividing by $p^3$ and adding \eqref{eq:moving-root-omitted} gives
\eqref{eq:advance:rank-second}, with the term $1$ of the bracket to spare
(it allows for the invariant of the diagonal case).

\emph{Step 6 (the first moment).} Let $U'=\mathbb A^1\setminus(S\cup A)$ and
$\mathcal G'=\mathcal F_S\otimes\mathcal L'$ on $U'$, with $\mathcal L'$ the Kummer
sheaf of $g_A$ on $U'$. $\mathcal G'$ is lisse of rank $h(s)\ge2$, pure of weight
$1$, tame, and geometrically irreducible (as in Step~2, by
\cref{lem:dl:sheaf}(i)--(iii) and \ref{it:prelim:E1}--\ref{it:prelim:E4}), so \cref{lem:tame-trace-bound} applies with $|S\cup A|$ finite punctures and
gives
$|\sum_{x\in U'(\Fp)}\mathrm{Tr}(\mathrm{Frob}_x\mid\mathcal G')|
\le h(s)(|S\cup A|-1)p\le(s+u-1)h(s)p$.
For $x\in U'(\Fp)$,
$\chi_A(x)\tau_S(x)=-p^{-1}\mathrm{Tr}(\mathrm{Frob}_x\mid\mathcal G')$; $\chi_A$
vanishes on $A$; and the at most $s$ points of $S\setminus A$ contribute at
most $p^{-1}\cdot s\cdot sr=s^2r^3$ to $\E_x\chi_A\tau_S$. So
$|\E_x\chi_A\tau_S|\le r^2(s+u-1)h(s)+r^3(s+u)s$, which is
\eqref{eq:advance:rank-first}.
\end{proof}

The sizes enter the bounds through the ranks $h(s)$, $h(t)$, the puncture bound $s+t+u$, and the term $r(s+t+u)st$ in \eqref{eq:advance:rank-second}. The root cuts use root sets of sizes $2$ to $11$ and
weights with at most ten roots, and they use the rank-sensitive bounds
themselves, through the pivot groups of \cref{subsec:advance:ordered-star}. The
product rows use root sets of sizes $2$ to $5$, weights with at most four
roots, and the uniform constant $512$ below; the constant $6000$ is the
corresponding uniform constant for the sizes of the root cuts.

\begin{corollary}[Uniform constants]\label{cor:moving-root-uniform}
For every odd prime $p$ and $S,T,A\subseteq\Fp$ with $2\le|S|,|T|\le11$ and
$|A|\le10$,
\[
  \Big|\E_x\chi_A\tau_S\tau_T-\frac{\one_{S=T,\ A=\emptyset}}p\Big|\le6000\,p^{-3/2}.
\]
If $|S|,|T|\le5$ and $|A|\le4$, the constant $6000$ can be replaced by $512$.
\end{corollary}

\begin{proof}
The bracket of \eqref{eq:advance:rank-second} is nondecreasing in $s$, $t$ and
$u$. At $(s,t,u)=(11,11,10)$ it is $31\cdot100+1+32\cdot121r=3101+3872r$
(rank $100$, at most $32$ finite punctures); at $(5,5,4)$ it is
$13\cdot16+1+14\cdot25r=209+350r$ (rank $16$, at most $14$ finite punctures).
Since $p\ge3$, $r\le1/\sqrt3$, and $3101+3872/\sqrt3<6000$,
$209+350/\sqrt3<512$.
\end{proof}

\begin{corollary}[Conditional $L^2$ bound]\label{cor:moving-root-conditional}
Let $y_1,\dots,y_d\in\Fp$ be distinct, $\varsigma_1,\dots,\varsigma_d\in\{\pm1\}$,
$E\subseteq\Fp$ a set containing $y_1,\dots,y_d$, and
$\Phi=\{x\in\Fp\setminus E:\varsigma_j(x-y_j)\in\Qsq\ \text{for}\ j\le d\}$; let
$\nu$ be the uniform law on $\Phi$, and suppose $|\Phi|\ge(1-e)p2^{-d}$ with
$e<1$, so that $\gamma=p2^{-d}/|\Phi|\le(1-e)^{-1}$. Let $S_1,\dots,S_m$ be
distinct subsets of $\Fp$ of sizes in $[2,\bar s]$ and $b_1,\dots,b_m\in\R$.
Then
\[
  \Big\|\sum_lb_l\tau_{S_l}\Big\|_{L^2(\nu)}
  \le\sqrt\gamma\,r\Big[\sum_lb_l^2+2^dC\,r\Big(\sum_l|b_l|\Big)^2\Big]^{1/2},
\]
with $C=6000$ if $\bar s\le11$ and $d\le10$, and $C=512$ if $\bar s\le5$ and
$d\le4$. In particular, for at most four pinned labels
$y_1,\dots,y_d$ ($d\le4$) and root sets of sizes $2$ to $5$,
the factor $2^dC$ may be taken equal to $8192$.
\end{corollary}

\begin{proof}
For $x\in\Phi$ each factor $1+\chi(\varsigma_j(x-y_j))$ equals $2$, and every
factor is nonnegative; so the density of $\nu$ relative to uniform probability is
at most $\gamma\prod_j\big(1+\chi(\varsigma_j(x-y_j))\big)$, and, the function
$F=\sum_lb_l\tau_{S_l}$ being real, $\E_\nu F^2\le\gamma\,\E_x\big[\prod_j(1+\chi(\varsigma_j(x-y_j)))F^2\big]$:
dropping the factor $\one[x\notin E]\le1$ from the nonnegative integrand only increases the right side. Expanding the product
over the subsets $P\subseteq\{1,\dots,d\}$, with $y(P)=\{y_j:j\in P\}$,
\[
  \E_x\Big[\prod_j\big(1+\chi(\varsigma_j(x-y_j))\big)F^2\Big]
  =\sum_P\chi\Big(\prod_{j\in P}\varsigma_j\Big)\sum_{l,l'}b_lb_{l'}
  \E_x\chi_{y(P)}\tau_{S_l}\tau_{S_{l'}}.
\]
By \cref{cor:moving-root-uniform} (arbitrary overlaps of $S_l$, $S_{l'}$ and the
pins are allowed), the main terms with $P=\emptyset$ and $l=l'$ give
$p^{-1}\sum_lb_l^2$, and all error terms, including the diagonal errors, together are at most
$2^dC\,r^3(\sum_l|b_l|)^2$ in modulus, since the root sets $S_l$ are distinct.
For $d\le4$ and $C=512$, $2^dC\le16\cdot512=8192$.
\end{proof}

\begin{remark}[One family of moments]\label{rem:trace-moment-family}
\Cref{thm:dl:H2} and \eqref{eq:advance:rank-first} are first moments, and
\eqref{eq:advance:rank-second} is a second moment, of the same normalized
traces: $\Sigma=p^2\E_z\chi_{Z(g)}\tau_f-[m\text{ odd}]\sum_z\chi_{Z(g)}(z)$ when
$f$ and $g$ are monic. Their common cohomological step is \cref{lem:tame-trace-bound}.
The output statements keep the data needed by their different applications: \cref{thm:dl:H2} keeps the exact number of common roots and the
symmetric variant, which the product rows and the completion variance use;
\cref{lem:moving-root-moments} keeps the ranks and the puncture counts, which
the pivot groups below need.
\end{remark}

\subsubsection{Grouping stars by a distinguished free label}\label{subsec:advance:ordered-star}

The conditional laws to which the moments are applied are those of a half
law given all labels but one.

\begin{definition}[Pinned half laws]\label{def:pinned-law}
Let $U$ be a finite set of $n$ vertices, and for each $i\in U$ let
$N^0(i)\subseteq U\setminus\{i\}$, $d_i=|N^0(i)|$ (the degree inside the half; $d^0_i$ in the notation of \cref{sec:spectral}), and signs
$\varsigma^0_{ij}\in\{\pm1\}$ ($j\in N^0(i)$). A probability law $\mu$ on the
injective maps $x:U\to\Fp$ is a \emph{pinned half law} with count errors
$e_i\in[0,1)$ if, for every $i\in U$ and every $x$ in the support of $\mu$, the
conditional law of $x_i$ given $x_{-i}=(x_j)_{j\ne i}$ is the uniform law on
\[
  \Phi_i(x_{-i})=\big\{y\in\Fp\setminus x(U\setminus\{i\}):\
  \varsigma^0_{ij}(y-x_j)\in\Qsq\ \text{for all}\ j\in N^0(i)\big\},
\]
and $|\Phi_i(x_{-i})|\ge(1-e_i)p2^{-d_i}$.
\end{definition}

For an oriented graph $G^+$ on $U$, the uniform law $\lambda_p(G^+)$ on
$\mathrm{Inj}_p(G^+)$ (\cref{def:graphs:labeling}) is a pinned half law, with
$N^0(i)$ the neighbours of $i$, $\varsigma^0_{ij}=+1$ for an edge $(j,i)$ and
$-1$ for an edge $(i,j)$: given $x_{-i}$, the admissible values of $x_i$ are
exactly $\Phi_i(x_{-i})$. By \cref{lem:extension}, with the at most $n-1$
other labels excluded, one may take $e_i=e_{d_i}(p,n)$ whenever this is below
one.

\begin{lemma}[Ordered almost orthogonality]\label{lem:ordered-orthogonality}
Let $x=(x_1,\dots,x_n)$ be a random vector with values in a finite set, with an
arbitrary joint law, and let $F=\sum_{i=1}^nF_i$, where each $F_i$ is a real
function of $(x_j)_{j\le i}$. Suppose that for every $i$, almost surely,
$\E[F_i^2\mid x_{-i}]\le\alpha_i^2$ and $|\E[F_i\mid x_{-i}]|\le\beta_i$. Then
\[
  \E F^2\le\sum_i\alpha_i^2+2\sum_i\beta_i\sum_{j<i}\alpha_j .
\]
\end{lemma}

\begin{proof}
$\E F^2=\sum_i\E F_i^2+2\sum_{j<i}\E F_iF_j$, and $\E F_i^2\le\alpha_i^2$. For
$j<i$ the function $F_j$ does not involve $x_i$, so it is a function of $x_{-i}$
and $\E F_iF_j=\E\big[F_j\,\E[F_i\mid x_{-i}]\big]$, of modulus at most
$\beta_i\E|F_j|\le\beta_i(\E F_j^2)^{1/2}\le\beta_i\alpha_j$.
\end{proof}

No independence between the $F_i$ is assumed: the cross terms are controlled
by conditioning on all variables but the later one, which turns the later
function into a combination of traces with fixed roots and the earlier one
into a constant.

\paragraph{Pivot groups.} Let $\mu$ be a pinned half law on $U$, fix a total
order of $U$, and let $a_T\in\R$ be coefficients indexed by the subsets
$T\subseteq U$ with $|T|\ge3$. Each such $T$ is assigned to its greatest vertex
$i(T)$, its \emph{pivot}; put $s_T=|T|-1$ and
$F_i(x)=\sum_{T:\,i(T)=i}a_T\,\tau_{x(T\setminus\{i\})}(x_i)$. For $i\in U$, with
all sums over the $T$ with $i(T)=i$ and with $d=d_i$, put
\begin{equation}\label{eq:pivot-sums}
\begin{gathered}
 A_i=\sum_Ta_T^2,\qquad L_i=\sum_T|a_T|,\qquad
 U_i=\sum_T|a_T|h(s_T),\\
 V_i=\sum_T|a_T|(s_T-1)h(s_T),\qquad
 W_i=\sum_T|a_T|s_T,\qquad X_i=\sum_T|a_T|s_T^2,\\
 E_i=2^d\Big[2U_iV_i+\big(\tfrac d2+1\big)U_i^2+L_i^2
   +r\big\{2W_iX_i+\tfrac d2W_i^2\big\}\Big],\qquad
 D_i^2=\frac{A_i+rE_i}{1-e_i},\\
 K_i=\frac{2^d}{1-e_i}\Big[V_i+\tfrac d2U_i
   +r\big\{X_i+\big(\tfrac d2+n-1\big)W_i\big\}\Big].
\end{gathered}
\end{equation}
These nonnegative expressions also apply when $d_i=0$, and they vanish for a
vertex that is the pivot of no support.

\begin{lemma}[Conditional moments of a pivot group]\label{lem:pivot-moments}
In this setting, for every $i\in U$ and every $x$ in the support of $\mu$,
\begin{enumerate}[label=\textup{(\alph*)}]
\item $\E_\mu[F_i^2\mid x_{-i}]\le r^2D_i^2$;
\item $|\E_\mu[F_i\mid x_{-i}]|\le r^2K_i$.
\end{enumerate}
\end{lemma}

\begin{proof}
Fix $i$ and $x_{-i}$, and write $d=d_i$, $X_{-i}=x(U\setminus\{i\})$ (a set of
$n-1$ labels), $\varsigma_j=\varsigma^0_{ij}$ and $y_j=x_j$ for $j\in N^0(i)$.
Given $x_{-i}$, $F_i=\sum_Ta_T\tau_{S_T}$ is a combination of traces with the
fixed root sets $S_T=x(T\setminus\{i\})$ of sizes $s_T\ge2$, and these root sets
are distinct, since $x$ is injective and distinct $T$ with the same pivot have
distinct sets $T\setminus\{i\}$. The conditional law of $x_i$ has density
$\gamma_i\prod_{j\in N^0(i)}\big(1+\chi(\varsigma_j(y-x_j))\big)\one[y\notin X_{-i}]$
relative to uniform probability on $\Fp$, with
$\gamma_i=p2^{-d}/|\Phi_i(x_{-i})|\le(1-e_i)^{-1}$. For $P\subseteq N^0(i)$ write
$x(P)=\{x_j:j\in P\}$ and $\varsigma_P=\prod_{j\in P}\varsigma_j$; then
$\prod_{j\in N^0(i)}(1+\chi(\varsigma_j(y-x_j)))=\sum_P\chi(\varsigma_P)\chi_{x(P)}(y)$.
Over the $2^d$ subsets, $\sum_P1=2^d$ and $\sum_P|P|=\frac d22^d$.

(a) As in the proof of \cref{cor:moving-root-conditional}, discarding the
restriction $y\notin X_{-i}$ from a square,
$\E_\mu[F_i^2\mid x_{-i}]\le\gamma_i\sum_P\chi(\varsigma_P)\sum_{T,T'}a_Ta_{T'}
\E_y\chi_{x(P)}\tau_{S_T}\tau_{S_{T'}}$. By \eqref{eq:advance:rank-second} with
$(s,t,u)=(s_T,s_{T'},|P|)$, the terms with $P=\emptyset$ and $T=T'$ give
$p^{-1}\sum_Ta_T^2=r^2A_i$, and the error is at most
\[
  r^3\sum_P\sum_{T,T'}|a_T||a_{T'}|\Big[(s_T+s_{T'}+|P|-1)h(s_T)h(s_{T'})+1
  +r(s_T+s_{T'}+|P|)s_Ts_{T'}\Big].
\]
Summing over $P$ replaces $|P|$ by $\frac d2$ and multiplies by $2^d$. Writing
$s_T+s_{T'}-1+\frac d2=(s_T-1)+(s_{T'}-1)+(1+\frac d2)$, the rank part sums to
$2U_iV_i+(\frac d2+1)U_i^2+L_i^2$ and the restored-point part to
$2W_iX_i+\frac d2W_i^2$. So the error is at most $r^3E_i$, and
$\E_\mu[F_i^2\mid x_{-i}]\le\gamma_i(r^2A_i+r^3E_i)\le r^2D_i^2$.

(b) Unlike the second moment, the first moment cannot drop the excluded
labels. Write $\one[y\notin X_{-i}]=1-\one[y\in X_{-i}]$:
\[
  \E_\mu[F_i\mid x_{-i}]=\gamma_i\sum_P\chi(\varsigma_P)\sum_Ta_T\E_y\chi_{x(P)}\tau_{S_T}
  -\frac{\gamma_i}p\sum_{y\in X_{-i}}\prod_{j\in N^0(i)}\big(1+\chi(\varsigma_j(y-x_j))\big)F_i(y).
\]
By \eqref{eq:advance:rank-first}, the first part is at most
$\gamma_ir^2\sum_P\sum_T|a_T|\big[(s_T+|P|-1)h(s_T)+r(s_T+|P|)s_T\big]
=\gamma_ir^22^d\big[V_i+\frac d2U_i+r(X_i+\frac d2W_i)\big]$ in modulus. In the
second part there are $n-1$ points, the pin product is at most $2^d$, and
$|F_i(y)|\le\sum_T|a_T|s_Tr=rW_i$ by \cref{lem:moving-root-trace}(b); so it is
at most $\gamma_ip^{-1}(n-1)2^drW_i=\gamma_ir^32^d(n-1)W_i$. This $(n-1)W_i$ term
covers the excluded labels exactly, and the sum is at most $r^2K_i$.
\end{proof}

Every vertex of a support with pivot $j$ precedes or equals $j$, so $F_j$ is a
function of $(x_k)_{k\le j}$, as \cref{lem:ordered-orthogonality} requires.

\begin{proposition}[Norm of a sum of higher star functions in a fixed vertex order]\label{prop:advance:ordered-star}
Let $\mu$ be a pinned half law on $U$, fix a total order of $U$ and real
coefficients $a_T$ ($T\subseteq U$, $|T|\ge3$), and let $F_i$, $D_i$, $K_i$ be as
in \eqref{eq:pivot-sums}. Then
\[
  \Big\|\sum_{i\in U}F_i\Big\|_{L^2(\mu)}\le M_\star
  :=r\Big[\sum_iD_i^2+2r\sum_iK_i\sum_{j<i}D_j\Big]^{1/2},
  \qquad |\E_\mu F_i|\le r^2K_i\quad(i\in U).
\]
In particular $\norm{P_0\sum_iF_i}_\mu\le M_\star$, where $P_0$ is the projection
off constants. No bound on $n$, on the degrees or on $|T|$ is needed; the root
cuts of the upper range have $n=36$, $d_i\le10$ and $|T|\le12$.
\end{proposition}

\begin{proof}
Apply \cref{lem:ordered-orthogonality}, with the vertices listed in the fixed
order, $\alpha_i=rD_i$ and $\beta_i=r^2K_i$ (\cref{lem:pivot-moments}(a),(b)):
$\E_\mu(\sum_iF_i)^2\le r^2\sum_iD_i^2+2r^3\sum_iK_i\sum_{j<i}D_j=M_\star^2$.
Integrating \cref{lem:pivot-moments}(b) gives the bound for $\E_\mu F_i$, and
$P_0$ is a contraction.
\end{proof}

\begin{remark}[The order, and monotonicity in $r$]\label{rem:pivot-order}
The order is only a choice of bound: the proposition holds for every total
order. The certificates place the vertices in decreasing internal degree,
breaking ties by their original indices, so every support uses one of its
vertices of smallest internal degree as its pivot. If each $e_i$ is replaced by
a bound $\bar e_i(r)<1$ that is a polynomial in $r$ with nonnegative
coefficients (as $e_t(p,n)$ of \cref{lem:extension} is), then $D_i$, $K_i$, and
hence $M_\star/r$ and the quantity $z_\star/r^2$ of \cref{lem:star-mean}, are
nondecreasing functions of $r$; so evaluation at the largest $r$ in the range (the smallest prime) bounds every smaller $r$, hence every larger prime (\cref{lem:majorant}).
\end{remark}

\paragraph{Diagonal and pivot bounds.}
The next lemma bounds the same kind of star marginal for a graph law, with
arbitrary real coefficients, by the $T=T'$ density bounds of part~\textup{(a)} and the conditional-mean bounds of part~\textup{(b)}
computed with \cref{thm:dl:H2} instead of conditional moments. It needs no
order of the vertices: each cross term is conditioned on all labels but that
of one vertex outside the other support.

\begin{lemma}[Weighted star marginals in $L^2$]\label{lem:wt:stars}
Let $G^+$ be an oriented graph on a set $U$ of $n$ vertices with degrees $d_i$
and maximum degree $\Delta^+$, let $p$ be an odd prime with
$e_{\Delta^+}(p,n)<1$, put $q=1/p$, and let $\mu=\lambda_p(G^+)$, a pinned half
law whose neighbour sets $N^0(i)$ are the sets of neighbours in $G^+$
(\cref{def:pinned-law}). For $T\subseteq U$ let $e(T)$ be the number of edges of
$G^+$ inside $T$ and $w_T(x)=\E_z\prod_{j\in T}\chi(z-x_j)$ as in
\cref{lem:star-trace}, and put
$R_U=\prod_{i\in U}\frac{1+e_{d_i}(p,n)}{1-e_{d_i}(p,n)}$ (the quantity $R_{\rm full}$ of \cref{lem:spec2:paley-one} when $U=U^+$). For $|T|\ge3$ and
$j\in T$ put $s_T=|T|-1$, $d=d_j$, $o=|N^0(j)\cap T|$,
\begin{gather}
 D^{\rm st}_T=\min\Big(2^{e(T)}R_U,\ \min_{i\in T}\frac{2^{d_i}}{1-e_{d_i}(p,n)},
 \ (|T|-1)^2\Big),\label{eq:wt:Dst}\\
 \beta_{T,j}=\frac q{1-e^-_d(n)}\Big[(d-o)+d(s_T-1)r
 +\sum_{a=2}^d\sum_{u\ge0}\binom ou\binom{d-o}{a-u}\mathcal B_{\min}(s_T,a,u)
 +2^d(n-1)s_Tr\Big],\label{eq:wt:beta}
\end{gather}
with the graded count error $e^-_d(n)$ of \eqref{eq:graphs:graded},
$\mathcal B_{\min}(m,a,u)=\min\big(\mathcal B(m,a,u),\mathcal B(a,m,u)\big)$ with
$\mathcal B$ of \eqref{eq:dl:B}, and binomial coefficients outside their range
zero. For $T\ne T'$ of size at least three let $\Xi_{TT'}$ be the least of the
numbers $\beta_{T,j}(D^{\rm st}_{T'})^{1/2}r$ ($j\in T\setminus T'$) and
$\beta_{T',j}(D^{\rm st}_T)^{1/2}r$ ($j\in T'\setminus T$). Then for all
$T,T'\subseteq U$ of size at least three:
\begin{enumerate}[label=\textup{(\alph*)}]
\item $\E_\mu w_T^2\le qD^{\rm st}_T$;
\item $|\E_\mu[w_T\mid x_{U\setminus\{j\}}]|\le\beta_{T,j}$ for $j\in T$ and
  every feasible $x_{U\setminus\{j\}}$;
\item $|\E_\mu w_Tw_{T'}|\le\Xi_{TT'}$ if $T\ne T'$.
\end{enumerate}
Consequently, for real coefficients $c_T$,
\[
 \Big\|\sum_{|T|\ge3}c_Tw_T\Big\|^2_\mu\le q\sum_Tc_T^2D^{\rm st}_T
 +\sum_{T\ne T'}|c_Tc_{T'}|\,\Xi_{TT'} .
\]
\end{lemma}

\begin{proof}
By \cref{lem:graphs:sequential}\ref{it:seq:count},\ref{it:seq:coarse},
$\mathrm{Inj}_p(G^+)\ne\emptyset$, so $\mu$ is defined, and it is supported on
injective labellings. Write $x_{-j}=x_{U\setminus\{j\}}$ and
$X_{-j}=\{x_k:k\ne j\}$. Since $e_t(p,n)$ is nondecreasing in $t$
(\ref{it:et:poly}), $e_{d_i}(p,n)\le e_{\Delta^+}(p,n)<1$ for every $i$, and
$e^-_d(n)\le e_d(p,n)<1$ (\ref{it:et:compare}).

(a) Any one of three bounds suffices. \emph{Weil.} On injective labellings
$\prod_{j\in T}(z-x_j)$ has $|T|$ distinct roots, so $|w_T|\le(|T|-1)r$ by
\cref{thm:weil}(a), and $\E_\mu w_T^2\le(|T|-1)^2q$. \emph{Density.} By
\cref{lem:graphs:vertex}\ref{it:vertex:tuple} with $V=T$, $C=\emptyset$ and any
ordering of $T$, the law of $x_T$ under $\mu$ has density at most
$2^{e(T)}R_{U\setminus T}\prod_{i\in T}(1-e_{t_i}(p,n))^{-1}$ with respect to the
uniform probability on $\Fp^T$, where $t_i\le d_i$ is the number of earlier
neighbours of $i$ in $T$, so that $\sum_it_i=e(T)$; as
$(1-e_{t_i})^{-1}\le(1-e_{d_i})^{-1}\le(1+e_{d_i})/(1-e_{d_i})$, the density is at
most $2^{e(T)}R_U$. Extend $w_T$ to all of $\Fp^T$ by its defining formula. By
\eqref{eq:two-root}, with $s=|T|\ge3$,
\begin{align*}
 \E_{v\in\Fp^T}w_T(v)^2&=\E_{z,z'}\prod_{j\in T}\E_{v_j}\chi(z-v_j)\chi(z'-v_j)
 =\E_{z,z'}(\one_{z=z'}-q)^s\\
 &=q\big[(1-q)^s+(p-1)(-q)^s\big]\le q(1-q),
\end{align*}
since $(1-q)^s\le(1-q)^2$ and $(p-1)q^s\le(p-1)q^2$; so
$\E_\mu w_T^2\le2^{e(T)}R_U\,q$. \emph{Conditional.} Fix $i\in T$, let
$f=\prod_{l\in T\setminus\{i\}}(X-x_l)$ and
$\mathcal G(\xi)=p^{-1}\sum_{y\in\Fp}\chi\big(f(y)(y-\xi)\big)$, so that
$w_T=\mathcal G(x_i)$. Since
$\sum_\xi\chi\big((y-\xi)(y'-\xi)\big)=p\one_{y=y'}-1$ by \eqref{eq:two-root},
\[
 \sum_{\xi\in\Fp}\mathcal G(\xi)^2
 =p^{-2}\sum_{y,y'}\chi\big(f(y)f(y')\big)\big(p\one_{y=y'}-1\big)
 =p^{-2}\Big[p(p-|T|+1)-\Big(\sum_y\chi(f(y))\Big)^2\Big]\le1 .
\]
By \cref{lem:graphs:vertex}\ref{it:vertex:dens}, given a feasible $x_{-i}$ the
label $x_i$ is uniform on a set $\Phi_i$ of at least
$p2^{-d_i}(1-e_{d_i}(p,n))$ elements, so
$\E_\mu[w_T^2\mid x_{-i}]\le|\Phi_i|^{-1}\sum_\xi\mathcal G(\xi)^2
\le q2^{d_i}/(1-e_{d_i}(p,n))$. Averaging over $x_{-i}$, and taking the least of
the three bounds and of the choices of $i$, gives (a).

(b) Fix $j\in T$ and a feasible $x_{-j}$; let $d$, $o$, $s_T$ be as in
\eqref{eq:wt:beta}, and $f$, $\mathcal G$ as in (a) with $i=j$, so that
$w_T=\mathcal G(x_j)$ and $\deg f=s_T\ge2$. Given $x_{-j}$, the label $x_j$ is
uniform on the admissible set $\Phi_j=\Phi_j(x_{-j})$ of \cref{def:pinned-law}
(\cref{lem:graphs:vertex}\ref{it:vertex:dens}). Put $\varsigma_k=\varsigma^0_{jk}$,
$\varsigma_P=\prod_{k\in P}\varsigma_k$ and $x(P)=\{x_k:k\in P\}$, and use the
character weights $\chi_A$ of \cref{subsec:moving-root-moments}. For
$\xi\notin X_{-j}$ every $\xi-x_k$ is nonzero, and
\[
 \one_{\Phi_j}(\xi)=\prod_{k\in N^0(j)}\tfrac12\big(1+\chi(\varsigma_k)\chi(\xi-x_k)\big)
 =2^{-d}\sum_{P\subseteq N^0(j)}\chi(\varsigma_P)\chi_{x(P)}(\xi).
\]
Hence, with $g_P=\prod_{k\in P}(X-x_k)$,
\[
 \sum_{\xi\in\Phi_j}\mathcal G(\xi)=2^{-d}\sum_{P\subseteq N^0(j)}\chi(\varsigma_P)
 \Big[\frac{\Sigma_P}p-\sum_{\xi\in X_{-j}}\chi_{x(P)}(\xi)\mathcal G(\xi)\Big],
 \qquad
 \Sigma_P=\sum_{y,\xi\in\Fp}\chi\big(f(y)\,g_P(\xi)\,(y-\xi)\big).
\]
\emph{$P=\emptyset$:} $\Sigma_\emptyset=\sum_y\chi(f(y))\sum_\xi\chi(y-\xi)=0$.
\emph{$P=\{k\}$:} $\sum_\xi\chi\big((\xi-x_k)(y-\xi)\big)=\chi(-1)(p\one_{y=x_k}-1)$
by \eqref{eq:two-root}, so $\Sigma_P=\chi(-1)\big[p\chi(f(x_k))-\sum_y\chi(f(y))\big]$
and $|\Sigma_P|/p\le|\chi(f(x_k))|+(s_T-1)r$ by \cref{thm:weil}(a); as $x$ is
injective, $\chi(f(x_k))=0$ exactly when $k\in T$. So the $d$ sets with $|P|=1$
contribute at most $(d-o)+d(s_T-1)r$. \emph{$|P|=a\ge2$:} $f$ and $g_P$ are monic
and split with distinct roots, of degrees $s_T\ge2$ and $a\ge2$, and by
injectivity they have exactly $u=|P\cap T|$ common roots; \cref{thm:dl:H2} with
$c=c'=1$, in both of its forms, gives $|\Sigma_P|\le p\,\mathcal B_{\min}(s_T,a,u)$,
and there are $\binom ou\binom{d-o}{a-u}$ such sets $P$ for given $a$ and $u$.
\emph{The excluded labels:} $|\mathcal G(\xi)|\le s_Tr$ for every $\xi$. If $\xi$
is not a root of $f$, $f(y)(y-\xi)$ has $s_T+1$ distinct roots and
\cref{thm:weil}(a) applies; if it is, $f(y)(y-\xi)=(y-\xi)^2\tilde f(y)$ with
$\tilde f$ of degree $s_T-1\ge1$ with distinct roots, the term $y=\xi$ vanishes,
and $|p\,\mathcal G(\xi)|\le|\sum_y\chi(\tilde f(y))|+1\le(s_T-2)\sqrt p+1\le
s_T\sqrt p$. The $2^d$ sets $P$ and the $n-1$ labels of $X_{-j}$ therefore
contribute at most $2^d(n-1)s_Tr$. Altogether
$|\sum_{\xi\in\Phi_j}\mathcal G(\xi)|$ is at most $2^{-d}$ times the bracket of
\eqref{eq:wt:beta}. Finally
$\E_\mu[w_T\mid x_{-j}]=|\Phi_j|^{-1}\sum_{\xi\in\Phi_j}\mathcal G(\xi)$, and
\cref{lem:extension}\ref{it:ext:lower}, with the $d$ distinct roots $x_k$
($k\in N^0(j)$), the signs $\chi(\varsigma_k)$ and the excluded set $X_{-j}$, which
has $n-1\le n$ elements and contains the roots, gives
$|\Phi_j|\ge p2^{-d}(1-e^-_d(n))>0$. Dividing gives (b).

(c) Let $j\in T\setminus T'$. Then $w_{T'}$ is a function of $x_{T'}$ and
$T'\subseteq U\setminus\{j\}$, so
$\E_\mu w_Tw_{T'}=\E_\mu\big[w_{T'}\,\E_\mu[w_T\mid x_{-j}]\big]$, and by (b) and (a)
$|\E_\mu w_Tw_{T'}|\le\beta_{T,j}\E_\mu|w_{T'}|\le\beta_{T,j}(qD^{\rm st}_{T'})^{1/2}
=\beta_{T,j}(D^{\rm st}_{T'})^{1/2}r$. The case $j\in T'\setminus T$ is symmetric,
and one of the two sets is nonempty as $T\ne T'$. For the last assertion expand
$\norm{\sum_Tc_Tw_T}^2_\mu=\sum_{T,T'}c_Tc_{T'}\E_\mu w_Tw_{T'}$, and bound the
diagonal terms by (a) and the others by (c).
\end{proof}

For a soft cross kernel the higher-star part of the marginal density of one half relative to $\mu$ is $\sum_{|T|\ge3}c_Tw_T$, where $T$ runs over sets of leaves in that half
and $c_T$ is the sum, over the centres in the other half, of the star
coefficients of support $T$ (\cref{cor:star-matrix}); for the cut kernel $K_0$ of \cref{sec:spectral} (all soft coefficients equal to $1$ on $E_\times$) $c_T=a_T$ of \eqref{eq:star-function}. The coefficients are arbitrary, so the
signs and sizes of the soft entries play no role in the lemma. The certificates
of the dense comparisons and of the six-query mixture
(\cref{prop:dense-comparison,prop:le-mixture}) evaluate the bound with the count
error $\big(\sum_{k\ge3}\binom dk(k-1)\big)r+\big(\binom d2+2^dn\big)q$ in place of
$e^-_d(n)$. Since $\sum_{k\ge3}\binom dk(k-1)=a_d+b_d-\binom d2$ (proof of
\cref{lem:extension}, and \ref{it:et:compare}), this error exceeds $e^-_d(n)$ by
$\big(b_d-\binom d2\big)(r-q)\ge0$ ($b_d=2^{d-1}-1\ge\binom d2$ for $d\ge1$); the
certificates use this weaker form: $\beta_{T,j}$ increases with $e^-_d(n)$ through the factor $(1-e^-_d(n))^{-1}$ in \eqref{eq:wt:beta}, so replacing $e^-_d(n)$ by a larger error below $1$ only enlarges the bound. For fixed coefficients, $D^{\rm st}_T$,
$\beta_{T,j}/q$ and $\Xi_{TT'}/(qr)$ are nondecreasing functions of $r$: they are
built from the polynomials $e_t(p,n)$, $e^-_t(n)$ and $\mathcal B(m,a,u)$ in $r$
with nonnegative coefficients (\ref{it:et:poly}, \eqref{eq:dl:B}) by sums,
products, minima and the maps $x\mapsto(1-x)^{-1}$, $x\mapsto(1+x)/(1-x)$ on
$[0,1)$. So the bound divided by $q$ is nondecreasing in $r$, and a check at one
prime covers every larger prime at which the coefficients are the same
(\cref{lem:majorant}).

\subsubsection{The complete-star mean}\label{subsec:advance:star-normalizer}

The constant $Z_0=\E_{\mu\otimes\mu}K_0$ normalizing the cut kernel $K_0$ of \cref{sec:spectral} needs the mean of the star function, not only
its centred part; here the affine constants of \cref{lem:star-trace} must be
kept.

\begin{lemma}[The complete-star mean]\label{lem:star-mean}
Consider spectral data with the star function $h_\star$ and the coefficients
$a_T$ of \eqref{eq:star-function}. Let $\mu=\lambda_p(G^+)$ be a pinned half law
with count errors $e_i$, fix a total order of $U^+$, let $F_i$, $K_i$, $M_\star$
be the quantities of \cref{subsec:advance:ordered-star} for the coefficients
$a_T$ with $|T|\ge3$, and put
\[
  B_{\rm even}=\sum_{T:\ |T|\ \text{even}\ \ge2}a_T
  =\sum_{k\in U^+}\Big(\frac12\Big[\prod_{j\in N(k)}\big(1+\chi(\varsigma_{jk})\big)
  +\prod_{j\in N(k)}\big(1-\chi(\varsigma_{jk})\big)\Big]-1\Big).
\]
Then $h_\star=\sum_iF_i-B_{\rm even}/p$ on $\mathrm{Inj}_p(G^+)$. Consequently
$h_\star-\E_\mu h_\star=P_0\sum_iF_i$, $\norm{h_\star-\E_\mu h_\star}_\mu\le M_\star$,
and
\[
  |\E_\mu h_\star|\le z_\star:=r^2\Big(\sum_iK_i+|B_{\rm even}|\Big).
\]
\end{lemma}

\begin{proof}
By \eqref{eq:star-function} and \cref{lem:star-trace}, the supports with $|T|=1$
contribute nothing, those with $|T|=2$ contribute $-a_T/p$, and those with
$|T|\ge3$ contribute $a_T\big(\tau_{x(T\setminus\{i(T)\})}(x_{i(T)})-\one_{|T|\ \rm
even}/p\big)$. Summing, $h_\star=\sum_iF_i-p^{-1}\big(\sum_{|T|=2}a_T
+\sum_{|T|\ge4\ \rm even}a_T\big)=\sum_iF_i-B_{\rm even}/p$. For the second
expression of $B_{\rm even}$: for a centre $k$ with signs
$c_j=\chi(\varsigma_{jk})$, $\sum_{T\subseteq N(k),\ |T|\ \rm even}\prod_{j\in T}c_j
=\frac12\big[\prod_j(1+c_j)+\prod_j(1-c_j)\big]$, and the empty support, which
contributes $1$, is subtracted. (For $|N(k)|\ge1$ a centre whose signs $c_j$ are
all equal contributes $2^{|N(k)|-1}-1$, and every other centre contributes
$-1$.) Then $\E_\mu h_\star=\sum_i\E_\mu F_i-B_{\rm even}/p$ and
$h_\star-\E_\mu h_\star=P_0\sum_iF_i$, and \cref{prop:advance:ordered-star} gives
both bounds, as $1/p=r^2$.
\end{proof}

\begin{remark}[How the mean enters, and what is not done]\label{rem:star-mean-scope}
The row stars have the same mean as the column stars, and together the stars
contribute $2\E_\mu h_\star$ to the normalizer of the cut kernel
(\cref{lem:spec:stars}(c)). This is how $z_\star$ and
$Z_\star\ge2|\E_\mu h_\star|$ enter hypothesis \ref{it:record:star-H3} of
\cref{thm:record:direct-star}, and how the same bound enters the normalizer
error of route~A (\cref{thm:spectral-criterion}). The number $z_\star$ depends
on the class of $p$ through the signs $\chi(\varsigma_{jk})$: the
root caps proved by Route~B, applied at $p\equiv3\pmod4$, use its value for that class
(\cref{lem:record:direct-star-inputs}(ii)), and the completed-root caps use the
larger of its two values (\cref{lem:upper-completed-root}). Two cautions
apply. The affine constants are kept: the two-leaf stars and the constants
$-\one_{|T|\ \rm even}/p$ make up $B_{\rm even}/p$, which does not vanish in the
mean. And the estimate bounds $\E_\mu h_\star$ itself; the bound for the centred part $P_0h_\star=h_\star-\E_\mu h_\star$ is not used in its place.

For one half law the mean is known exactly. On the half reference of the middle
root cut L20, the graph law on the six vertices $(a_0,a_1,b_0,b_1,c_0,c_1)$ whose only edges are the four arrows $a_i\to c_j$ ($i,j\in\{0,1\}$) from the pair $(a_0,a_1)$ to the pair $(c_0,c_1)$, the remaining two vertices $(b_0,b_1)$ being isolated, the two-leaf star means
cancel, the three-leaf ones add up to zero at each centre, and the four-leaf
ones are given
exactly in terms of the graph moment $M_{G_5}$ of the variance credits
(\cref{lem:l20-star-mean}, proved in \cref{sec:variance} next to that moment).
The normalizer of that cut uses this exact value instead of the general bound
(\cref{subsec:middle-l20}).
\end{remark}

\subsection{Where these estimates enter}\label{subsec:second-order-uses}

\Cref{tab:second-order-uses} lists every use in the final construction. The
sheaves of \cref{lem:dl:sheaf} stand behind all rows except those of
\cref{thm:dl:K,cor:dl:normalized,cor:raw-pair-maps}, which rest on the
representation theory of $\mathrm{PGL}_2(\Fp)$, and the last row, an external
evaluation through a modular form (\cref{subsec:external:modular}).

\begin{table}[ht]
\small
\caption{Where the second-order estimates enter.}\label{tab:second-order-uses}
\begin{tabular}{@{}>{\raggedright\arraybackslash}p{3.1cm}%
>{\raggedright\arraybackslash}p{5.6cm}>{\raggedright\arraybackslash}p{5.4cm}@{}}
\toprule
estimate & enters as & in\\
\midrule
\cref{thm:dl:K}, \cref{cor:dl:normalized}(a)
 & the norm $2r^2$ of a two-vertex strongly connected block of a witness kernel
   (\cref{lem:matching:kernel}), hence the complete-rectangle coefficient of the
   compressed partitions (\cref{lem:advance:compressed-partition})
 & the nonstar bounds of every spectral cap: the twelve-label grid, the caps of
   the root cut L20 (\cref{subsec:middle-l20,sec:positive-roots}), the ten
   nonroot upper cuts, $N_{\rm ns}$ of the direct-star root
   (\cref{lem:record:direct-star-inputs}(iii)), the pinned rows of \cref{subsec:push:pinnedremainder}\\
\cref{cor:dl:normalized}
 & the injective $K_{2,2}$ core $2r^2$ of a matched neighbourhood in the
   completion variance (\cref{lem:push:localvariance}, through
   \cref{lem:matching:kernel}(b))
 & the coupling block (\cref{cor:spec2:coupling}) of the completed-root caps (\cref{lem:upper-completed-root})
   from $2^{52}$\\
\cref{cor:raw-pair-maps}
 & the remainder $\norm R^2\le c_0q^2$ of the raw pair maps (b), and the norms
   $r$ and $\sqrt{q(1-q)}$ of the maps of (a)
 & the square identity, the linear trace and the layer squares of the
   positive roots
   (\cref{lem:root-square,lem:root-linear,prop:root-layer})\\
\cref{thm:dl:H2}
 & the pair term of the completion variance, through the coefficient triples of
   $\mathcal B(a,b,0)$ (\cref{rem:dl:H2-sharp}(ii))
 & the coupling block of the completed-root caps (\cref{lem:upper-completed-root}) from $2^{52}$\\
\cref{thm:dl:H2}
 & the complete conditional means (PR1) of the product rows, with
   $p\,\mathcal B(4,4,0)<40p$ (\cref{rem:dl:H2-sharp}(iii),
   \cref{subsec:advance:product-rows})
 & the row inputs of the upper bands\\
\cref{lem:wt:stars} (through \cref{thm:dl:H2})
 & the $L^2$ norm of the weighted higher-star marginals of a soft cross kernel
 & the dense comparisons and the six-query mixture of the upper range
   (\cref{prop:dense-comparison,prop:le-mixture})\\
\cref{cor:moving-root-uniform,cor:moving-root-conditional}
 & the weighted moments (PR7) and the conditional $L^2$ bounds (PR8) of the
   product rows
 & the row inputs of the upper bands\\
\cref{prop:advance:ordered-star,lem:star-mean}
 & $M_\star$, $z_\star$ and $Z_\star$ (\cref{lem:record:direct-star-inputs}(i),(ii))
 & the direct-star root on the upper bands from $P_3$ to $2^{52}$
   (\cref{thm:record:direct-star}) and the completed-root caps from $2^{52}$
   (\cref{lem:upper-completed-root})\\
\cref{thm:ahlgren-ono} (\cref{app:external})
 & the bounds $-(1-q)(2qr+q^2)\le M_{G_5}\le(1-q)(2qr-q^2)$ of the graph
   moment \eqref{eq:record:modular-moment}, in the second moment of a
   neighbourhood count (\cref{lem:record:weighted-moment}) and in the quartic
   star mean (\cref{lem:l20-star-mean})
 & the root-half variances $\chi_p$ (the $\chi^2$-divergences of \cref{sec:variance}, not the character $\chi_p$ above) of the middle roots
   (\cref{lem:record:root-marginal}) in the restorations and the completion of
   the root cut L20 (\cref{subsec:middle-l20}); the star mean of that cut; the
   linear prices and the normalization of the positive roots
   (\cref{lem:root-linear,lem:root-normalization})\\
\bottomrule
\end{tabular}
\end{table}

\paragraph{Consequences and what to read next.}
This section used the following inputs; its outputs are the three second-order estimates of the opening, whose uses are listed in \cref{tab:second-order-uses}.
From Part~II: Weil's bound (\cref{thm:weil}), Jacobi sums (\cref{thm:jacobi}),
the two-root identity \eqref{eq:two-root}, the Paley operator
(\cref{lem:square-shift}(c)), and, in \cref{rem:pivot-order} and in the discussion after \cref{lem:wt:stars} only (not in any proof), the admissible majorants of
\cref{lem:majorant}. From Part~IV: strict-square labellings and the extension
count (\cref{def:graphs:labeling,lem:extension}), to see that half laws are
pinned and, with the one-vertex conditional laws and tuple densities of graph
laws (\cref{lem:graphs:vertex}\ref{it:vertex:dens},\ref{it:vertex:tuple}), the
sequential counts (\cref{lem:graphs:sequential}), the graded lower count
\cref{lem:extension}\ref{it:ext:lower} and the properties
\ref{it:et:compare}, \ref{it:et:poly} of the count errors, in
\cref{lem:wt:stars}; spectral data
(\cref{def:spec:data}) for the cross signs of the star function; and, in remarks and discussion only (not in any proof), \cref{lem:spec:stars,cor:star-matrix,lem:swapped-rectangle}. From
\cref{app:external}: the character table of $\mathrm{GL}_2(\Fp)$ with its
class sizes (\cref{thm:prelim:gl2}); the $\ell$-adic facts
\ref{it:prelim:E1}--\ref{it:prelim:E4}; the trace formula, the cohomology of
open curves, Deligne's theorem and the Euler characteristic formula
(\cref{thm:prelim:lefschetz,thm:prelim:curves,thm:prelim:weights,thm:prelim:gos});
and the Katz--Sarnak family with its self-duality and its local monodromy at the
finite roots (\cref{thm:prelim:ks}).
